% Modern editorial, markup and figure/layout contributions: CC-BY-SA-4.0.
% Historical text and illustrations: Leonhard Euler (1744), public domain.
% Component scope and upstream attribution: see RIGHTS.md.
% Chapter V, original pp. 171--227. Collated against scan.
\subsection*{Caput V.}
\textit{Methodus, inter omnes curvas eadem proprietate præditas, inveniendi eam, quæ maximi minimive proprietate gaudeat.}
\subsection*{Definitio.}
\originalsection{1}\sourcepage{171}Proprietas communis est Formula integralis, seu expressio indefinita, quæ in omnes curvas ex quibus quæsitam determinari oportet, æqualiter competit.

\subsection*{Scholion I.}
\originalsection{2}Hactenus Methodum maximorum ac minimorum tradidimus absolutam, in qua perpetuo, inter omnes omnino curvas eidem abscissæ respondentes, una requiri solebat, quæ maximi minimive cujuspiam proprietate gauderet. Nunc autem progredimur ad Methodum relativam, in qua unam lineam maximi minimive proprietate præditam determinare docebimus, non ex omnibus omnino lineis eidem abscissæ respondentibus, verum ex illis, innumerabilibus quidem, lineis curvis tantum, quibus una quædam proprietas proposita pluresve sint communes. Ac primo quidem, in hoc Capite, innumerabiles curvas eidem abscissæ respondentes contemplabimur, quæ unam quandam proprietatem habeant communem; ex hisque unam lineam investigabimus, in qua expressio quæcunque indefinita maximum minimumve obtineat valorem. Hoc in genere inprimis celebre est Problema Isoperimetricum, initio hujus sæculi publice propositum, in quo, inter omnes curvas ejusdem longitudinis quæ quidem eidem abscissæ respondeant, eam definiri oportebat, quæ contineret maximi minimive cujuspiam proprietatem. Postmodum autem hæc Quæstio in latiori sensu est accepta, ut ista deter\sourcepage{172}minatio non solum inter omnes curvas ejusdem longitudinis fieret, verum etiam inter omnes curvas alia quacunque proprietate communi præditas; quam ipsam quæstionem in hoc Capite pertractare suscepimus. Cum igitur curva sit eligenda, non ex omnibus omnino curvis eidem abscissæ respondentibus, verum ex iis, innumerabilibus duntaxat, in quas proprietas quæpiam proposita æqualiter competat; hanc ipsam proprietatem ante omnia considerari oportet, quam hic nomine proprietatis communis indicamus. Hæc igitur proprietas communis, veluti æqualitas longitudinis curvarum, omnia puncta media afficere debet, \& hanc ob rem erit functio indefinita, quæ, non ex unici curvæ elementi, verum ex totius curvæ positione determinetur. Quam ob rem istiusmodi proprietas communis erit, vel formula integralis indefinita simplex, vel expressio plures ejusmodi formulas integrales complectens. Omnino igitur pari modo erit comparata, quo ipsa maximi minimive formula, seu expressio. Eædem igitur varietates atque divisiones, quas ante circa maximi minimive expressionem fecimus \& tractavimus, æque ad proprietatem communem pertinebunt.

\subsection*{Corollarium I.}
\originalsection{3}Si igitur proprietas communis fuerit proposita, quæ sit $B$, tum omnes curvæ sunt considerandæ, quæ pro eadem data abscissa eundem valorem ipsius $B$ continent; atque ex his ea debet definiri, quæ habeat maximum vel minimum.

\subsection*{Corollarium II.}
\originalsection{4}In Problematis ergo huc pertinentibus duas res datas esse oportet, proprietatem communem $B$, ac maximi minimive expressionem $A$. Quibus datis, inter omnes curvas pro data abscissa eundem valorem $B$ continentes, ea definiri debebit, quæ pro eadem abscissa valorem ipsius $A$ habeat maximum vel minimum.

\subsection*{Corollarium III.}
\originalsection{5}\sourcepage{173}Dantur autem non solum infinitæ curvæ, quæ pro data abscissa eandem proprietatem communem habeant, sed etiam dantur infinitis modis. Assumta enim curva quacunque pro lubitu, ea determinatum habebit valorem proprietatis communis propositæ; præter eam autem dabuntur innumerabiles aliæ eundem valorem proprietatis communis pro eadem abscissa continentes.

\subsection*{Corollarium IV.}
\originalsection{6}Proposita igitur expressione quacunque indefinita, innumerabilia infinitarum curvarum dabuntur genera; quorum quodlibet genus infinitas in se complectitur curvas, quæ pro eadem data abscissa eundem illius expressionis valorem contineant.

\subsection*{Corollarium V.}
\originalsection{7}Cum igitur infinita dentur genera, quorum singula innumerabiles lineas curvas comprehendant, in quas proposita pro proprietate communi expressio æqualiter competat; in uno quoque genere dabitur una curva, quæ, pro reliquis ejusdem generis curvis, alteram expressionem in maximo minimove gradu contineat.

\subsection*{Corollarium VI.}
\originalsection{8}Quoniam ergo, ex quolibet genere, una curva maximi minimive proprietate prædita invenitur; omnino ejusmodi curvæ satisfacientes infinitæ invenientur, quarum quævis ita erit comparata, ut inter omnes alias eadem proprietate communi gaudentes, maximi minimive proprietate sit prædita.

\subsection*{Scholion II.}
\originalsection{9}\sourcepage{174}Hæc omnia magis illustrabuntur, si proprietatem communem, de qua hactenus in genere sumus locuti, definiamus. Sit igitur proprietas communis, formula longitudinem arcus curvæ exprimens, maximi minimive expressio autem sit $\int Zdx$; ita ut, inter omnes curvas quæ habeant arcus eidem abscissæ respondentes inter se æquales, ea debeat determinari, in qua pro eadem abscissa fiat $\int Zdx$ maximum vel minimum. Manifestum autem est, non solum infinitas lineas curvas dari pro eadem abscissa longitudine æquales, verum hoc etiam infinitis modis fieri posse. Sit enim abscissa communis $=a$, sumaturque quæcunque longitudo $c$ major quam $a$, infinitæ exhiberi poterunt lineæ, tum rectæ tum curvæ, quarum singularum longitudo sit $=c$; atque inter has definiri poterit una, in qua sit $\int Zdx$ maximum vel minimum. Loco $c$ autem infinitæ accipi possunt quantitates; eo quod alia non adest conditio, nisi ut sit $c>a$; atque quilibet valor pro $c$ assumtus dabit unam curvam maximi minimive proprietate præditam. Quamobrem, pro infinitis ipsius $c$ valoribus, infinitæ reperientur lineæ curvæ quæstioni satisfacientes. Neque tamen idcirco Quæstio pro indeterminata est habenda: nam solutio ipsa, infinitas curvas satisfacientes præbens, ita est interpretanda, ut unaquæque harum curvarum inventarum inter omnes alias æque longas possideat valorem formulæ $\int Zdx$ in maximo minimove gradu. Perspicuum autem est, quod hic de æqualibus arcubus curvarum ostendimus, idem de alia quacunque formula seu expressione indeterminata valere debere. Ita si, inter omnes curvas quæ, pro data abscissa $x=a$, valorem formulæ $\int Ydx$ eundem continent, ea requiratur in qua sit $\int Zdx$ maximum vel minimum; tum infinitæ quidem reperientur lineæ satisfacientes: verum hæ inter se ita discrepabunt, ut quælibet, inter omnes alias possibiles lineas curvas secum valorem formulæ $\int Ydx$ communem habentes, contineat formulæ $\int Zdx$ valorem maximum vel minimum.

\subsection*{Propositio I. Theorema.}
\originalsection{10}\sourcepage{175}Quæ curva, inter omnes omnino curvas eidem abscissæ respondentes, maximi minimive cujuspiam propositi proprietate gaudet; eadem curva simul, inter omnes curvas communi quacunque cum ipsa proprietate præditas, eadem maximi minimive proprietate gaudebit.

\subsection*{Demonstratio.}
Sit maximi minimive expressio $=A$, proprietas autem communis $=B$; eritque tam $A$ quam $B$, vel formula integralis indefinita, vel expressio ex hujusmodi pluribus formulis composita. Ponamus jam curvam esse inventam, quæ, inter omnes omnino curvas eidem abscissæ respondentes, expressionem $A$ contineat maximam vel minimam; ea curva certum quendam expressionis $B$ continebit valorem; præter eam autem dabuntur innumerabiles aliæ, in quas idem expressionis $B$ valor competet; hæque innumerabiles curvæ omnes jam continentur in illis omnibus omnino curvis, ex quibus ea, in qua expressio $A$ est maximum minimumve, est inventa. Cum igitur hæc curva, inter omnes omnino curvas, proposita maximi minimive proprietate gaudeat; eadem quoque, inter illas infinitas curvas secum expressionem $B$ communem habentes, valorem expressionis $A$ maximum minimumve possidebit. Q. E. D.

\subsection*{Corollarium I.}
\originalsection{11}Methodus igitur absoluta etiam Problematis Methodi relativæ resolvendis inservit: dum unam semper curvam satisfacientem exhibet. Verum tamen solutionem completam non largitur.

\subsection*{Corollarium II.}
\originalsection{12}Curva ergo, quæ, inter omnes, expressionem $A$ habet\sourcepage{176} maximam vel minimam, erit una ex infinitis illis curvis, quarum singulæ, inter omnes alias secum communi proprietate $B$ gaudentes, eandem expressionem $A$ maximam habent minimamve.

\subsection*{Corollarium III.}
\originalsection{13}Solutio igitur Problematis, quo, inter omnes curvas eadem communi proprietate $B$ præditas, ea quæritur in qua sit $A$ maximum vel minimum, latius patebit, quam si absolute, inter omnes curvas, ea quæreretur in qua est $A$ maximum vel minimum; illaque solutio hanc tanquam casum specialem in se comprehendet.

\subsection*{Propositio II. Problema.}
\originalsection{14}Methodum resolvendi Problemata, in quibus, inter omnes curvas communi quadam proprietate gaudentes, ea requiritur quæ maximi minimive cujuspiam propositi proprietate gaudeat, in genere adumbrare.

\subsection*{Solutio.}
\marginnote{\footnotesize \hyperlink{fig:15}{Fig.~15.}}Omne maximum vel minimum ita est comparatum, ut, facta mutatione infinite parva, valor ejus omnino non immutetur. Quamobrem si curva $az$, inter omnes curvas eidem abscissæ $AZ$ respondentes, quæ quidem communi proprietate $B$ gaudeant, habeat valorem expressionis $A$ maximum vel minimum; eundem valorem retinebit, si ipsi talis mutatio infinite parva inferatur, qua communis proprietas $B$ non turbetur. Ad hoc autem non sufficit, ut ante fecimus, unicam applicatam, puta $Nn$, particula infinite parva $nv$ auxisse: quoniam enim hoc modo tota mutatio unica conditione determinatur, per eam effici nequit, ut tam proprietas communis $B$ in ipsam curvam \& immutatam æqualiter competat, quam maximi minimive expressio $A$. Quocirca mutationem adhibendam binis conditio\sourcepage{177}nibus
 determinatum esse oportet; id quod obtinebitur, si binæ applicatæ $Nn$, \& $Oo$ particulis infinite parvis $nv$ \& $o\omega$ augeantur. Quod si ergo curva hoc modo immutari concipiatur; primum efficiendum est, ut proprietas communis cum in ipsam curvam tum in mutatam æque competat; deinde etiam maximi minimive expressio in utraque curva eundem valorem retinere debebit. Prius præstabitur, si expressionis, qua proprietas communis continetur, valor differentialis investigetur, oriundus ex translatione binorum $n$ \& $o$ in $v$ \& $\omega$, isque evanescens ponatur: posteriori vero conditioni satisfiet, si pari modo valor differentialis expressionis, quæ maximum minimumve esse debet, quæratur, oriundus ex binis particulis $nv$ \& $o\omega$, atque nihilo æqualis ponatur. Hoc pacto, duæ obtinebuntur æquationes, altera ex proprietate communi, altera ex maximi minimive expressione; utraque autem ejusmodi habebit formam $S.nv+T.o\omega=0$; in qua $S$ \& $T$ erunt quantitates ad curvam pertinentes. Ex binis autem ejusmodi æquationibus eliminabuntur particulæ $nv$ \& $o\omega$; provenietque æquatio pro curva quæsita, quæ, inter omnes alias eadem communi proprietate $B$ præditas, habeat valorem expressionis $A$ maximum vel minimum. Q. E. I.

\subsection*{Corollarium I.}
\originalsection{15}Solutio igitur hujusmodi Problematum quoque reducitur ad inventionem valorum differentialium: ipsi autem valores differentiales ab iis quos ante dedimus in hoc discrepant, quod ex translatione duorum curvæ punctorum definiri debeant.

\subsection*{Corollarium II.}
\originalsection{16}Ejusmodi valores differentiales ergo ex duabus particulis $nv$ \& $o\omega$ oriundos, in quovis Problemate, binos investigari oportet; alterum pro proprietate communi, alterum pro maximi minimive expressione.

\subsection*{Corollarium III.}
\originalsection{17}\sourcepage{178}Inventis autem in quovis Problemate his duobus valoribus differentialibus, uterque nihilo æqualis poni debet; ex quo binæ nascentur æquationes, quæ, eliminandis particulis assumtis $nv$ \& $o\omega$, præbebunt unam æquationem naturam curvæ quæsitæ exprimentem.

\subsection*{Corollarium IV.}
\originalsection{18}Si ergo, inter omnes curvas eidem abscissæ respondentes, quæ communi proprietate $B$ æqualiter sunt præditæ, ea requiratur, in qua expressio $A$ fiat maximum vel minimum; tum utriusque expressionis $A$ \& $B$ valores differentiales, ex binis particulis $nv$ \& $o\omega$ oriundi, quæri, \& nihilo æquales poni debent; ex quibus duabus æquationibus si eliminentur particulæ $nv$ \& $o\omega$, emerget æquatio pro curva quæsita.

\subsection*{Corollarium V.}
\originalsection{19}In hac itaque operatione, ambæ expressiones $A$ \& $B$ omnino pariter tractantur; neque in considerationem venit, utra vel proprietatem communem vel maximum minimumve denotet. Ex quo perspicuum est, eandem solutionem prodire debere, si expressiones $A$ \& $B$ inter se commutentur.

\subsection*{Corollarium VI.}
\originalsection{20}Eadem ergo solutio locum habebit, sive, inter omnes curvas communi proprietate $B$ gaudentes, ea quæratur in qua sit $A$ maximum vel minimum: sive vicissim, inter omnes curvas communi proprietate $A$ gaudentes, ea quæratur in qua sit $B$ maximum vel minimum.

\subsection*{Scholion.}
\originalsection{21}\sourcepage{179}Ambas expressiones $A$ \& $B$, licet in se spectatæ res omnino diversas significent, inter se commutabiles esse ipsa solutionis natura sponte patet. Quod si enim ad binas particulas $nv$ \& $o\omega$ respiciamus, quibus applicatæ $Nn$ \& $Oo$ augentur; primum eas ita comparatas esse oportet, ut proprietas communis $B$, tam in ipsa curva quam in mutata, eundem valorem obtineat; scilicet proprietas communis $B$ in curvam $amnopz$ \& in $amv\omega pz$ æque competere debet: deinde pari modo per easdem particulas $nv$ \& $o\omega$ efficiendum est, ut expressio $A$, quæ maximum minimumve esse debet, tam pro curva $amnopz$ quam pro $amv\omega pz$ eundem valorem recipiat. Atque adeo, tam proprietas communis, quam maximi minimive natura, eandem plane conditionem in calculum inducit; ex quo manifestum est ambas expressiones datas, quarum altera proprietatem communem, altera maximi minimive rationem continet, inter se commutari atque confundi posse, salva Solutione. Hanc ob rem ergo, in Solutione hujusmodi Problematum, sufficit nosse ambas illas expressiones; neque ad Solutionem absolvendam nosse opus est, utra proprietatem communem aut maximum minimumve significet. Sic si, inter omnes curvas longitudine æquales, quæratur ea, quæ maximam aream comprehendat; eadem reperitur curva quæ prodit, si, inter omnes curvas æquales areas includentes, ea quæratur quæ sit brevissima, vel minimam longitudinem habeat. Hæc ita se habent, si maximi minimive quod quæritur natura ita fuerit comparata, ut ejus valor differentialis sit $=0$. Jam supra autem animadvertimus, duplicis generis dari maxima \& minima, in quorum altero valor differentialis sit $=0$, in altero vero $=\infty$. Hic vero tantum maxima ac minima prioris generis contemplamur: nam, in hac Methodo relativa, posterius genus locum omnino habere nequit. Quod si enim valor differentialis, qui convenit maximi minimive expressioni, infinite magnus ponatur; tum ex hoc solo æquatio pro curva reperitur; neque ideo proprietas communis in computum ingreditur. Quare, si hujus
\sourcepage{180}generis maximum vel minimum in Methodo absoluta locum habet, eadem curva in Methodo relativa eadem proprietate gaudebit, quæcunque proprietas communis adjungatur. Cum igitur totum Solutionis hujusmodi Problematum momentum versetur in inventione valorum differentialium, qui ex binis particulis $nv$ \& $o\omega$ oriuntur; Methodum trademus, ejusmodi valores differentiales pro quacunque expressione indeterminata inveniendi, eo modo, quo supra usi sumus ad inveniendos valores differentiales ex unica particula $nv$ oriundos.

\subsection*{Propositio III. Problema.}
\originalsection{22}\marginnote{\footnotesize \hyperlink{fig:15}{Fig.~15.}}Proposita quacunque expressione indeterminata, quæ ad datam abscissam $AZ$ referatur; invenire ejus valorem differentialem, ortum ex translatione binorum curvæ punctorum $n$ \& $o$ in $v$ \& $\omega$.

\subsection*{Solutio.}
Ponamus abscissam $AI=x$, \& applicatam $Ii=y$, erit $Kk=y'$, $Ll=y''$, $Mm=y'''$, $Nn=y^{\mathrm{iv}}$, $Oo=y^{\mathrm{v}}$, $Pp=y^{\mathrm{vi}}$ \&c. Harum applicatarum duæ tantum, nempe $y^{\mathrm{iv}}$ \& $y^{\mathrm{v}}$ patiuntur alterationem a particulis $nv$ \& $o\omega$ ipsis adjunctis. Erit igitur applicatæ $y^{\mathrm{iv}}$ valor differentialis $=nv$, \& applicatæ $y^{\mathrm{v}}$ valor differentialis $=o\omega$, reliquarum vero applicatarum omnium valor differentialis erit $=0$. Hinc reliquarum quantitatum ad curvam pertinentium $p$, $q$, $r$, $s$, \&c. valores differentiales habebuntur, quatenus eæ ab his binis applicatis $y^{\mathrm{iv}}$ \& $y^{\mathrm{v}}$ pendent. Sic cum sit $p=\frac{y'-y}{dx}$, erit valor differentialis ipsius $p=0$; similiterque ipsius $p'$, \& $p''$: at cum sit $p'''=\frac{y^{\mathrm{iv}}-y'''}{dx}$, erit ipsius $p'''$ valor differentialis $=\frac{nv}{dx}$; \&, ob $p^{\mathrm{iv}}=\frac{y^{\mathrm{v}}-y^{\mathrm{iv}}}{dx}$, erit valor differentialis ipsius $p^{\mathrm{iv}}=\frac{o\omega}{dx}-\frac{nv}{dx}$; porroque ipsius $p^{\mathrm{v}}$ erit $=-\frac{o\omega}{dx}$. Deinde cum sit
\sourcepage{181}$q=\frac{p'-p}{dx}$; erit valor differentialis ipsius $q''=\frac{nv}{dx^2}$; ipsius $q'''=\frac{o\omega}{dx^2}-\frac{2nv}{dx^2}$, ipsius $q^{\mathrm{iv}}=-\frac{2o\omega}{dx^2}+\frac{nv}{dx^2}$, ipsius $q^{\mathrm{v}}=\frac{o\omega}{dx^2}$. Hocque modo similiter progredi licet ad sequentes quantitates $r$, $s$, \&c. cum suis derivativis; hincque nascetur sequens Tabella, qua singularum harum quantitatum valores differentiales exhibentur.
\[
\renewcommand{\arraystretch}{1.7}
\begin{array}{l|l}
d.y^{\mathrm{iv}}=nv&d.q''=\dfrac{nv}{dx^2}\\
d.y^{\mathrm{v}}=o\omega&d.q'''=-\dfrac{2nv}{dx^2}+\dfrac{o\omega}{dx^2}\\
d.p'''=\dfrac{nv}{dx}&d.q^{\mathrm{iv}}=\dfrac{nv}{dx^2}-\dfrac{2o\omega}{dx^2}\\
d.p^{\mathrm{iv}}=-\dfrac{nv}{dx}+\dfrac{o\omega}{dx}&d.q^{\mathrm{v}}=\dfrac{o\omega}{dx^2}\\
d.p^{\mathrm{v}}=-\dfrac{o\omega}{dx}&
\end{array}
\]
\[
\renewcommand{\arraystretch}{1.7}
\begin{array}{l|l}
d.r'=+\dfrac{nv}{dx^2}&d.s=+\dfrac{nv}{dx^4}\\
d.r''=-\dfrac{3nv}{dx^3}+\dfrac{o\omega}{dx^3}&d.s'=-\dfrac{4nv}{dx^4}+\dfrac{o\omega}{dx^4}\\
d.r'''=+\dfrac{3nv}{dx^3}-\dfrac{3o\omega}{dx^3}&d.s''=+\dfrac{6nv}{dx^4}-\dfrac{4o\omega}{dx^4}\\
d.r^{\mathrm{iv}}=-\dfrac{nv}{dx^3}+\dfrac{3o\omega}{dx^3}&d.s'''=-\dfrac{4nv}{dx^4}+\dfrac{6o\omega}{dx^4}\\
d.r^{\mathrm{v}}=-\dfrac{o\omega}{dx^3}&d.s^{\mathrm{iv}}=+\dfrac{nv}{dx^4}-\dfrac{4o\omega}{dx^4}\\
&d.s^{\mathrm{v}}=+\dfrac{o\omega}{dx^4}
\end{array}\qquad\&c.
\]
Ex hac Tabella perspicitur, in valoribus differentialibus totidem terminos particula $o\omega$ affectos occurrere, ac particula $nv$; atque in utrisque pares adesse coefficientes: discrimen vero in hoc consistere, ut cuilibet termino particula $o\omega$ affecto respon\sourcepage{182}deat quantitas immediate sequens eam, cui respondet similis terminus particula $nv$ affectus. Sic dum terminus $-\frac{2nv}{dx^2}$ reperitur in valore differentiali quantitatis $q'''$, ita terminus $-\frac{2o\omega}{dx^2}$ adest in valore differentiali quantitatis sequentis $q^{\mathrm{iv}}$. Deinceps, ob duplicis generis terminos in valoribus differentialibus occurrentes, quorum alteri particulam $nv$, alteri particulam $o\omega$ involvunt, valor differentialis cujuscunque expressionis indeterminatæ hujusmodi habebit formam $nv.I+o\omega.K$; de qua primum, manifestum est membrum prius $nv.I$ esse ejusdem expressionis valorem differentialem, qui oritur si sola particula $nv$ consideretur; eritque ideo $nv.I$ ille ipse valor differentialis, quem supra pro quavis expressione oblata definire docuimus; ita ut hoc membrum per præcepta supra tradita pro quavis expressione indeterminata exhibere liceat. Quod ad alterum membrum $o\omega.K$ attinet, quia singuli termini in quibus $o\omega$ inest perpetuo respondent quantitatibus sequentibus eas, quibus respondent similes termini particulam $nv$ involventes, palam est quantitatem $K$ fore valorem, quem quantitas $I$ in proximo sequente loco induit, atque idcirco esse $K=I'=I+dI$. Quare cum membrum $nv.I$ ex præceptis jam supra datis assignare queamus, ex eo porro alterum membrum $o\omega.K=o\omega(I+dI)$ innotescet. Sit igitur $V$ expressio quæcunque indeterminata, cujus valorem differentialem ex duabus particulis $nv$ \& $o\omega$ oriundum definiri oporteat. Ponamus ejus valorem differentialem ex unica particula $nv$ oriundum esse $=nv.I$; eritque valor differentialis, qui ex ambabus particulis $nv$ \& $o\omega$ oritur, $=nv.I+o\omega.I'$, seu $=nv.I+o\omega(I+dI)$; qui igitur ope regularum supra datarum facile assignari poterit. Q. E. I.

\subsection*{Corollarium I.}
\originalsection{23}Omnium ergo expressionum, quarum valores differentiales ex unica particula $nv$ oriundos invenire docuimus, earundem
\sourcepage{183}valores differentiales ex binis particulis $nv$ \& $o\omega$ oriundos nunc definire possumus.

\subsection*{Corollarium II.}
\originalsection{24}Hæc igitur Methodus valebit tam ad expressionum valores differentiales inveniendos, qui non pendent a quantitate abscissæ propositæ $AZ$, quam qui ab istius abscissæ longitudine pendent.

\subsection*{Corollarium III.}
\originalsection{25}Quin etiam si expressio proposita, quæ vel proprietatem communem continet, vel maximum minimumve esse debet, fuerit functio duarum pluriumve formularum integralium; ejus valor differentialis ex binis particulis $nv$ \& $o\omega$ oriundus eadem lege definietur.

\subsection*{Scholion.}
\originalsection{26}In Capitibus superioribus vidimus valorem differentialem cujuscunque expressionis, qui ex unica particula $nv$ oritur, hujusmodi habere formam $nv.dx.T$, seu $nv.Tdx$; ubi $T$ denotat quantitatem finitam: quare ejusdem expressionis valor differentialis ex binis particulis $nv$ \& $o\omega$ ortus erit $=nv.Tdx+o\omega.T'dx$, quemadmodum in Solutione ostendimus. Eadem autem forma facile ad hunc modum potest evinci: Scilicet si ponatur $o\omega=0$, tum prodire debet ipse valor differentialis ex unica particula $nv$ ortus; quem supra invenire docuimus, eritque $nv.Tdx$. Sin autem ponatur $nv=0$, ac sola particula $o\omega$ consideratur, valor differentialis simili modo reperietur quo supra usi sumus: non autem erit $=o\omega.Tdx$; nam quia particula $o\omega$ in situ sequente accipitur, loco $T$ ejus valor sequens pariter sumi debet; ita ut valor differentialis verus futurus sit $=o\omega.T'dx$. Quod si ergo utraque particula $nv$ \& $o\omega$ conjunctim consideretur, erit valor differentialis $=nv.Tdx+o\omega.T'dx$; eo quod in ipso cal\sourcepage{184}culo particulæ $nv$ \& $o\omega$ nusquam inter se permiscentur, sed utraque perpetuo seorsim tractari possit. Ut autem hæc ad notandi modum in superiori capite receptum accommodemus; ponamus $V$ esse expressionem quamcunque indeterminatam, quæ, pro abscissa definita $AZ=a$, valorem recipiat $=A$; ejusque valorem differentialem ex particula $nv$ ortum esse $=nv.dA$; ubi $dA$ nobis denotet idem quod ante $Tdx$; poteritque iste valor $dA$ ex expressione $V$, modo in Capitibus præcedentibus exposito, inveniri. Hoc invento erit ejusdem expressionis $V$ valor differentialis ex binis particulis $nv$ \& $o\omega$ oriundus $=nv.dA+o\omega.dA'$ ubi $dA'$ denotat valorem $dA$ suo differentiali auctum. Quanquam autem ista valorum differentialium ex binis particulis oriundorum ad nostrum institutum omnino est necessaria; tamen solutio ipsa Problematum huc pertinentium eo iterum reducetur, ut per solos valores differentiales modo supra exposito inventos, qui scilicet ex unica particula $nv$ nascuntur, absolvi queat; id quod ex sequente Propositione mox patebit.

\subsection*{Propositio IV. Problema.}
\originalsection{27}Inter omnes curvas ad eandem datam abscissam $AZ=a$ relatas, in quas idem valor expressionis indefinitæ $W$ competit; determinare eam, in qua sit expressio $V$ maximum vel minimum.

\subsection*{Solutio.}
Ponamus curvam $az$ quæsito satisfacere, atque expressionem $W$ in ea obtinere valorem determinatum $=B$; erit ergo hæc curva $az$ inter omnes alias curvas ad eandem abscissam $AZ$ relatas, in quibus expressio $W$ eundem obtinet valorem, ita comparata, ut in ea expressio $V$ maximum minimumve valorem recipiat, qui sit $=A$. Ad curvam ergo hanc inveniendam, positis abscissa indefinita $AI=x$, \& applicata respondente $Ii=y$; binæ applicatæ $Nn$ \& $Oo$ particulis infinite parvis $nv$ \& $o\omega$ augeri concipiantur: quo facto, tam ipsius $W$ quam ipsius $V$ va\sourcepage{185}lor differentialis, qui ex his duabus particulis $nv$ \& $o\omega$ adjunctis nascetur nihilo æqualis poni debebit, uti in Propositione secunda ostendimus. Sit jam expressionis $V$ valor differentialis ex unica particula $nv$ ortus $=nv.dA$, atque expressionis alterius $W$ valor differentialis ex eadem unica particula $nv$ ortus $=nv.dB$, quos valores differentiales ex præceptis in superioribus Capitibus datis invenire licebit. Nunc igitur, dum binas particulas $nv$ \& $o\omega$ contemplamur, erit expressionis $V$ valor differentialis $=nv.dA+o\omega.dA'$; alterius vero expressionis $W$ valor differentialis erit $=nv.dB+o\omega.dB'$. Quocirca, ad quæsitam curvam inveniendam, fieri oportet cum $nv.dA+o\omega.dA'=0$, tum etiam $nv.dB+o\omega.dB'=0$. Multiplicentur ambæ æquationes per quantitates quascunque, ita ut prodeat
\[
\begin{aligned}
nv.\alpha dA+o\omega.\alpha dA'&=0\\
nv.\beta dB+o\omega.\beta dB'&=0.
\end{aligned}
\]
Fiatque ad particulas $nv$ \& $o\omega$ eliminandas tam $\alpha dA+\beta dB=0$ quam $\alpha dA'+\beta dB'=0$; eruntque $\alpha$ \& $\beta$ ejusmodi quantitates, sive constantes sive variabiles, quæ utrique æquationi satisfaciunt. Quoniam vero est $\alpha dA+\beta dB=0$, erit quoque $\alpha'dA'+\beta'dB'=0$; quæ æquatio, cum $\alpha dA'+\beta dB'=0$ comparata, monstrat esse debere $\alpha'=\alpha$, \& $\beta'=\beta$; ex quo quantitates hæ $\alpha$ \& $\beta$ debunt esse constantes, \& quidem quæcunque. Sumtis itaque pro $\alpha$ \& $\beta$ quantitatibus quibuscunque constantibus, æquatio pro curva erit $\alpha dA+\beta dB=0$. Hæc eadem æquatio prodit, si methodo consueta particulas $nv$ \& $o\omega$ eliminemus. Erit nempe
\[
\frac{nv}{o\omega}=-\frac{dA'}{dA}=-\frac{dB'}{dB},\qquad
\frac{dA'}{dA}=\frac{dB'}{dB},\quad\text{seu}\quad\frac{ddA}{dA}=\frac{ddB}{dB},
\]
ob $dA'=dA+ddA$, \& $dB'=dB+ddB$. Æquatio autem $\frac{ddA}{dA}=\frac{ddB}{dB}$ integrata dat $ldA=ldB+lC$. Seu $dA=
CdB$;\sourcepage{186} quæ, posito $C=-\frac{\beta}{\alpha}$, transit in $\alpha dA+\beta dB=0$; quam ipsam ante invenimus. Quamobrem ad Problema resolvendum oportet, tam expressionis proprietatem communem continentis $W$, quam expressionis quæ maximum minimumve esse debet $V$, valores differentiales, methodo in superioribus Capitibus tradita, investigare, eosque per quantitates constantes quascunque multiplicare, summamque $=0$ ponere; quo facto, resultabit æquatio naturam curvæ quæsitæ exprimens. Q. E. I.

\subsection*{Corollarium I.}
\originalsection{28}Nunc igitur, ad Quæstiones in hac Propositione contentas resolvendas, sufficit nosse valores differentiales ex unica particula $nv$ oriundos; quos supra jam expedite invenire docuimus.

\subsection*{Corollarium II.}
\originalsection{29}Quare ad hoc negotium in subsidium vocari debebit Caput præcedens IV, ex eoque cum §.~7 tum §.~31. In loco priore enim continentur præcepta valores differentiales inveniendi, si expressiones indeterminatæ propositæ fuerint formulæ integrales singularis, in altero vero, si sint functiones duarum pluriumve ejusmodi formularum integralium.

\subsection*{Corollarium III.}
\originalsection{30}Proposita ergo proprietate communi $W$, \& maximi minimive expressione $V$, utriusque expressionis valorem differentialem ex his præceptis quæri oportet: quibus inventis, \& per constantes arbitrarias multiplicatis, eorum aggregatum nihilo æquale positum dabit æquationem pro curva quæsita.

\subsection*{Corollarium IV.}
\originalsection{31}Si, inter omnes omnino curvas eidem abscissæ $AZ$ res\sourcepage{187}pondentes, quæratur ea, in qua expressio $V$ maximum minimumve obtineat valorem; pro ea habetur ista æquatio $dA=0$; denotante $dA$ valorem differentialem expressionis $V$.

\subsection*{Corollarium V.}
\originalsection{32}Quod si autem, inter omnes curvas eidem abscissæ $AZ$ respondentes, in quas expressio $W$ æqualiter competat, quæratur ea in qua expressio $V$ maximum minimumve habeat valorem; invenitur pro ea ista æquatio $\alpha dA+\beta dB=0$.

\subsection*{Corollarium VI.}
\originalsection{33}Perspicuum ergo est, curvam, quæ, inter omnes omnino curvas, habeat $V$ maximum vel minimum, cujus æquatio est $dA=0$, contineri in æquatione $\alpha dA+\beta dB=0$, qua exprimitur curva, quæ, inter omnes eadem communi proprietate $W$ gaudentes, habeat $V$ maximum vel minimum.

\subsection*{Corollarium VII.}
\originalsection{34}In ipsa igitur prima æquatione, quam Solutio præbet, $\alpha dA+\beta dB=0$; jam inest una constans arbitraria; quæ autem per id determinari debet, ut valor expressionis $W$ datum obtineat valorem.

\subsection*{Corollarium VIII.}
\originalsection{35}Problema itaque sic solvi poterit, ut, inter omnes curvas eidem abscissæ $AZ$ respondentes, in quibus expressio $W$ eundem datum obtineat valorem, definiatur ea in qua sit valor ipsius $V$ maximus vel minimus.

\subsection*{Corollarium IX.}
\originalsection{36}Ex his denique intelligitur, Solutionem Problematis pro\sourcepage{188}positi, convenire cum Solutione hujus Problematis, quo, inter omnes omnino curvas eidem abscissæ $AZ$ respondentes, requiratur ea quæ habeat $\alpha V+\beta W$ maximum vel minimum. Quæ quæstio, etsi ad Methodum absolutam pertineat, tamen dat æquationem $\alpha dA+\beta dB=0$, quam ipsam invenimus.

\subsection*{Scholion I.}
\originalsection{37}Ex his igitur non solum Methodus facilis atque expedita colligitur, omnes Quæstiones huc pertinentes resolvendi; verum etiam natura hujus generis Problematum penitius cognoscitur. Primo enim apparet, quod jam supra demonstravimus, Solutionem eandem fore, sive, inter omnes curvas communi proprietate $W$ præditas, quæratur ea quæ habeat $V$ maximum vel minimum; sive inverse, inter omnes curvas communi proprietate $V$ præditas, ea requiratur in qua sit $W$ maximum vel minimum. Deinde etiam intelligitur, Quæstionem ita proponi posse, ut ejus Solutio ad Methodum maximorum ac minimorum absolutam pertineat; congruit enim Problema propositum cum hoc, quo, inter omnes omnino curvas ad eandem abscissam $AZ$ relatas, requiritur ea in qua sit ista expressio $\alpha V+\beta W$ maximum vel minimum; atque hæc Problematis transformatio in causa est, quod Solutio per valores differentiales ex unica particula $nv$ oriundos perfici queat, neque amplius opus sit duas hujusmodi particulas considerare, prouti primo intuitu natura Quæstionis postulare videbatur. Hanc autem convenientiam postmodum, per se, ac sine ista Methodo qua binæ particulæ considerantur, demonstrabimus; quo veritas hæc, summi in isto negotio momenti, magis confirmetur. Ad solvendas cæterum hujusmodi Quæstiones, ante oculos habere oportet præcepta Capite præcedente in compendium redacta; quorum ope valores differentiales quarumcunque expressionum inveniri poterunt. Primo enim, §.~7 illius Capitis recensentur Casus, quibus formularum integralium solitariarum valores exhibentur: tum vero §.~31 traditur Methodus inveniendi valores differentiales expressionum\sourcepage{189}, quæ ex duabus pluribusve formulis integralibus utcunque sint compositæ. Ex his itaque subsidiis, pro quavis Quæstione oblata, tam maximi minimive expressionis quam proprietatis communis valor differentialis assignari poterit: utroque autem invento, æquatio pro Curva quæsita nullo negotio formabitur; cum tantum opus sit aggregatum quorumcunque multiplorum illorum binorum valorum differentialium nihilo æquale poni. Hæcque æquatio inventa, deinceps pari modo erit tractanda, quo supra, cum in reductione ad construendum, tum in integratione usi sumus.

\subsection*{Scholion II.}
\originalsection{38}Jam observavimus in æquatione $\alpha dA+\beta dB=0$, quam Solutio immediate suppeditat, unam inesse quantitatem constantem; quæ autem non omnino sit arbitraria, sed ex conditione proposita debeat determinari. Scilicet, cum in omnes curvas ex quibus quæsitam definiri oportet, eadem expressio $W$ æqualiter competere debeat, seu in omnibus eundem valorem, puta $B$, obtinere; hæc quantitas $B$ tanquam data spectari potest; atque cum ipsa in calculum non ingrediatur, ita constantes $\alpha$ \& $\beta$ definire licebit, ut valor expressionis $W$, abscissæ $AZ=a$ respondens, ipsi $B$ æqualis fiat; hocque pacto, Quæstio alioquin indeterminata determinabitur. Eatenus autem tantum determinabitur, quatenus, per integrationes post instituendas, novæ constantes arbitrariæ etiam per totidem puncta definiuntur. Prorsus nimirum ut ante, totidem puncta præscribi poterunt, per quæ curva quæsita transeat, quot novæ constantes per integrationes ingredi censendæ sunt. Horum autem numerus innotescet ex gradu differentialium summo, qui in æquatione inerit. Quoniam vero tota Quæstio ad Methodum absolutam revocari potest, numerus istiusmodi constantium perpetuo erit par; seu æquatio resultans $\alpha dA+\beta dB=0$, erit vel finita, vel differentialis secundi gradus, vel differentialis quarti gradus, vel differentialis sexti gradus, vel octavi, vel ita porro. Quod
\sourcepage{190}si æquatio prodit finita, tum quoque curva penitus jam erit determinata; siquidem ratio inter $\alpha$ \& $\beta$ ita definiatur, ut expressio $W$ datum recipiat valorem $B$ in curva inventa, quam determinationem perpetuo adhiberi ponimus. Si æquatio inveniatur differentialis secundi gradus, tum duobus punctis curva inventa determinabitur; congruum autem ac more receptum est ipsos curvæ terminos $a$ \& $z$ præscribi, hisque casibus Problema determinabitur, si conditio ista adjungatur, ut curva quæsita intra datos terminos $a$ \& $z$ contineatur. Sin autem æquatio prodeat differentialis quarti gradus, tum quatuor punctis pro lubitu assignatis, curva satisfaciens determinabitur; hæc igitur definiri ita conveniet, ut, præter terminos extremos $a$ \& $z$, simul positio tangentium in his terminis præscribatur. Sin perveniatur ad æquationem differentialem sexti gradus, tum curva per sex quæcunque puncta determinabitur: eorum autem loco præscribi poterunt primo ambo termini $a$ \& $z$, tum positio tangentium in his terminis, ac tertio curvedo in his ipsis locis seu radii osculi quantitas. His igitur notatis, intelligetur ex ipsa Solutione cujusmodi conditio ad Problematis cujusque propositionem adjungi debeat, ut id fiat penitus determinatum: hæcque admonitio, non solum hic, sed etiam in Methodo absoluta atque reliqua Methodo relativa, locum habet.

\subsection*{Scholion III.}
\originalsection{39}Discrimen hic etiam maximi momenti inprimis est notandum, ex quo in Methodo absoluta primariam tractationis partitionem desumsimus. Consistit id autem in modo, quo curva inventa Quæstioni satisfacit. Fieri enim potest, ut ejus quæcunque portio ad abscissam indefinitam relata requisita proprietate gaudeat; deinde etiam dantur casus, quibus nonnisi ea portio quæ definitæ abscissæ $AZ=a$ respondet, conditioni Problematis satisfaciat. Illud scilicet evenit, si quantitas hæc $a$ in æquationem, quam Solutio suppeditat, vel omnino non ingreditur, vel in quantitates arbitrarias $\alpha$ \& $\beta$ comprehendi
\sourcepage{191}queat. Ex quo manifestum est, si ambæ formulæ $W$ \& $V$ in casu primo §.~7 Capitis præcedentis recensito contineantur, tum curvæ inventæ quamlibet portionem ad Quæstionem esse accommodatam. Deinde vero etiam fieri potest, ut licet quantitas $a$, seu quantitates ab ea pendentes, vel in alterutro valore differentiali insint, vel in utroque; tamen eæ vel se mutuo tollant in æquatione $\alpha dA+\beta dB=0$, vel sub arbitrariis $\alpha$ \& $\beta$ comprehendi queant; quo casu pariter quamvis curvæ inventæ portionem satisfacere oportet. Hoc autem tantum locum habet, si non datus ac determinatus præscribatur valor, quem proprietas communis $W$ in portione satisfaciente obtinere debeat: tum enim fieri nequit, ut in quavis portione eundem valorem sortiatur. Ex Solutione autem unius cujusque Quæstionis facile intelligetur, qua conditione, sive tota curva $az$, sive quævis portio, satisfacere queat; id quod commodissime in Exemplis ostendi poterit.

\subsection*{Exemplum I.}
\originalsection{40}Inter omnes curvas ad abscissam $AZ$ relatas, in quibus formula $\int yxdx$ eundem obtinet valorem, invenire eam in qua sit valor formulæ $\int yydx$ minimus.

Erit igitur proprietas communis $W=\int xydx$, cujus, ob $d.xy=ydx+xdy$, valor differentialis est $=nv.dx.x$. Maximi autem minimive formula est $V=\int yydx$, cujus, ob $d.yy=2ydy$, valor differentialis est $=nv.dx.2y$. Obtinebitur ergo, divisione per $nv.dx$ instituta, hæc æquatio $\alpha x+2\beta y=0$; ex qua patet Quæstioni satisfacere lineam rectam in $A$ cum axe $AZ$ angulum quemcunque constituentem. Et quia longitudo abscissæ $AZ=a$ non in computum ingreditur, quævis hujus rectæ portio æque satisfaciet. Quod si autem postuletur, ut pro data abscissa $AZ=a$, formula $\int yxdx$ datum obtineat valorem, puta $B$; tum ob $y=mx$, fiet $\int yxdx=\frac13 mx^3$; ideoque $\frac13 ma^3=B$; ex quo positio lineæ rectæ
\sourcepage{192}ita definietur, ut esse debeat $y=\frac{3Bx}{a^3}$. Hæc igitur recta jam ista proprietate gaudebit, ut ea inter omnes lineas, sive rectas, sive curvas, quæ pro data abscissa $AZ=a$ habeant formulæ $\int xydx$ valorem $=B$, producat formulæ $\int yydx$ minimum valorem.

\subsection*{Exemplum II.}
\originalsection{41}Inter omnes curvas ejusdem longitudinis puncta $a$ \& $z$ jungentes, invenire eam quæ maximam vel minimam aream $aAZz$ comprehendat.

Quoniam proprietas communis est longitudo arcus $=\int dx\sqrt{1+pp}$; erit ejus valor differentialis $=-nv.d.\frac{p}{\sqrt{1+pp}}$. Deinde maximi minimive formula est $\int ydx$, cujus valor differentialis est $nv.dx$: unde pro curva quæsita ista habebitur æquatio
\[
dx=bd.\frac{p}{\sqrt{1+pp}}.
\]
\& integrando
\[
x+c=\frac{bp}{\sqrt{1+pp}};
\]
ideoque $p=\frac{x+c}{\sqrt{b^2-(x+c)^2}}=\frac{dy}{dx}$. Hinc ergo integrando fit $y=f\pm\sqrt{b^2-(x+c)^2}$, seu $b^2=(y-f)^2+(x+c)^2$; quæ est æquatio generalis pro Circulo. Quamobrem arcus Circuli quicunque per puncta $a$ \& $z$ ductus, inter omnes alias lineas curvas ejusdem longitudinis, vel maximam vel minimam aream $aAZz$ includet. Duplici autem modo Circuli arcus datæ longitudinis intra terminos $a$ \& $z$ constitui potest; altero, quo concavitatem axi $AZ$ obvertit; altero, quo convexitatem. Priori casu manifestum est aream fore maximam, posteriore vero minimam. Atque hinc si dentur termini $a$ \& $z$, una cum longitudine curvæ intra hos terminos constitutæ, quam majorem quidem esse oportet lineam rectam hos terminos jungentem; Solutio penitus erit determinata: arcus Circuli enim hujus longitudinis per hos terminos poterit describi
\sourcepage{193}unicus, qui, prout vel concavitatem vel convexitatem axi $AZ$ obvertat, aream formabit vel maximam vel minimam.

\subsection*{Corollarium.}
\originalsection{42}\marginnote{\footnotesize \hyperlink{fig:16}{Fig.~16.}}Hinc etiam patet arcum circularem $az$, per terminos $a$ \& $z$ ductum, non solum maximam aream $aAZz$, inter omnes alias lineas ejusdem longitudinis formare; sed etiam quæcunque linea $aCEDz$ a termino $a$ ad terminum $z$ ducta detur, cum ea arcus circularis $az$ maximam includet aream. Nam si area $aAZz$ est maxima, erit quoque area $aAZz-aAC-zZD+CED$, ob areas $aAC$, $zZD$ \& $CED$ constantis magnitudinis quæcunque linea pro $az$ capiatur, maxima.

\subsection*{Exemplum III.}
\originalsection{43}\marginnote{\footnotesize \hyperlink{fig:7}{Fig.~7.}}Inter omnes curvas ejusdem longitudinis puncta $A$ \& $M$ jungentes, invenire eam quæ, cum rectis $AC$ \& $MC$ ad punctum fixum $C$ ductis, maximam vel minimam comprehendat aream $ACM$.

Quoniam, ob data puncta $A$, $C$, $M$, rectæ $AC$ \& $MC$ positione dantur, ponatur angulus $ACM=x$, seu descripto, centro $C$, radio $CB=1$, arcu circulari $BS$, sit hic arcus $BS=x$, \& ponatur $CM=y$; erit $Ss=dx$, $Mn=ydx$, \& area $ACM=\frac12\int yydx$. Porro ob $mn=dy$, erit $Mm=\sqrt{y^2dx^2+dy^2}=dx\sqrt{yy+pp}$, posito $dy=pdx$. Quare inter omnes æquationes relationem ipsarum $x$ \& $y$ continentes, quæ, pro dato ipsius $x$ valore, eandem præbent quantitatem $\int dx\sqrt{yy+pp}$, eam definiri oportet, quæ, pro eodem ipsius $x$ valore, præbeat formulæ $\frac12\int yydx$ quantitatem vel maximam vel minimam. Cum igitur formulæ $\int dx\sqrt{yy+pp}$ valor differentialis sit
\[
=nv.dx\left(\frac{y}{\sqrt{yy+pp}}-\frac1{dx}d.\frac{p}{\sqrt{yy+pp}}\right)
\]
\& formulæ $\int yydx$ valor differentialis $=\sourcepage{194}nv.dx.y$; habebitur pro curva quæsita ista æquatio
\[
ydx=\frac{bydx}{\sqrt{yy+pp}}-bd.\frac{p}{\sqrt{yy+pp}},
\]
quæ, per $p$ multiplicata, abit in hanc
\[
\begin{aligned}
ydy&=\frac{bydy}{\sqrt{yy+pp}}-bpd.\frac{p}{\sqrt{yy+pp}}\\
&=bd.\sqrt{yy+pp}-\frac{bpdp}{\sqrt{yy+pp}}-bpd.\frac{p}{\sqrt{yy+pp}};
\end{aligned}
\]
cujus integrale est
\[
\frac12 yy=b\sqrt{yy+pp}-\frac{bpp}{\sqrt{yy+pp}}+bc=\frac{byy}{\sqrt{yy+pp}}+bc.
\]
Ducatur in tangentem $MP$ ex $C$ perpendiculum $CP=u$; erit $u=\frac{yy}{\sqrt{yy+pp}}$, habebiturque $yy=2bu+bc$: quam æquationem supra jam ostendimus esse ad Circulum. Quamobrem arcus Circuli per terminos $A$ \& $M$ ductus hanc habebit proprietatem, ut, inter omnes alias curvas ejusdem secum longitudinis terminos $A$ \& $M$ jungentes, aream $ACM$ exhibeat vel maximam vel minimam; prout ille arcus, vel concavitatem, vel convexitatem intra angulum $ACM$ vertat. Quo ipso id confirmatur, quod §.~præced. in genere adnotavimus.

\subsection*{Exemplum IV.}
\originalsection{44}\marginnote{\footnotesize \hyperlink{fig:15}{Fig.~15.}}Inter omnes curvas puncta $a$ \& $z$ jungentes, quæ circa axem $AZ$ rotatæ generant solida ejusdem superficiei; determinare eam quæ simul producat volumen solidi hoc modo generati maximum.

Superficies solidi hoc modo generati, proportionalis invenitur formulæ integrali huic $\int ydx\sqrt{1+pp}$, cujus valor differentialis est
\[
nv.dx\left(\sqrt{1+pp}-\frac1{dx}d.\frac{yp}{\sqrt{1+pp}}\right).
\]
Volumen vero solidi hoc modo generati est ut $\int yydx$, cujus valor differentialis est $=nv.dx.2y$. Quocirca resultabit ista æquatio
\[2ydx=bdx\sqrt{1+pp}-bd.\frac{yp}{\sqrt{1+pp}}.\]
\sourcepage{195}Multiplicetur hæc per $p$, ut prodeat
\[
\begin{aligned}
2ydy&=bdy\sqrt{1+pp}-bpd.\frac{yp}{\sqrt{1+pp}}\\
&=bd.y\sqrt{1+pp}-\frac{bypdp}{\sqrt{1+pp}}-bpd.\frac{yp}{\sqrt{+pp}},
\end{aligned}
\]
cujus integrale est
\[yy=by\sqrt{1+pp}-\frac{bypp}{\sqrt{1+pp}}-bc=\frac{by}{\sqrt{1+pp}}+bc.\]
Erit ergo $by=(yy-bc)\sqrt{1+pp}$, \& $p=\frac{\sqrt{b^2y^2-(yy-bc)^2}}{yy-bc}=\frac{dy}{dx}$. Quare erit
\[dx=\frac{(yy-bc)dy}{\sqrt{bbyy-(yy-bc)^2}}.\]
De hac æquatione primo notandum est, si fuerit $c=0$ fore $dx=\frac{ydy}{\sqrt{bb-yy}}$; ideoque curvam esse Circulum, cujus centrum in axe $AZ$ sit positum; ille igitur arcus circularis, centro in axe $AN$ sumpto descriptus \& per data duo puncta $a$ \& $z$ transiens Quæstioni satisfaciet; erit autem is unicus, ideoque solidum definitæ superficiei generabit. Quare si, inter omnes curvas quæ solida alius atque diversæ superficiei generant, quæratur ea quæ maximum volumen producat, ea non erit Circulus, sed alia curva in æquatione $dx=\frac{(yy-bc)dy}{\sqrt{bbyy-(yy-bc)^2}}$ contenta. Non solum enim, ob binas constantes $b$ \& $c$, effici potest, ut curva per præscripta duo puncta $a$ \& $z$ transeat; sed etiam ut longitudo curvæ $az$ existat datæ magnitudinis. Cæterum longitudo curvæ, ob $\int dx\sqrt{1+pp}=\int\frac{bydx}{yy-bc}$ fiet
\[
=\int\frac{bydy}{\sqrt{bbyy-(yy-bc)^2}},
\]
cujus integrale a quadraturæ Circuli pendet, estque
\[
=\frac b2\operatorname{A\,cos\textnormal{.}}\frac{b(2c+b)-2yy}{b\sqrt{bb+4bc}}+\textit{Const.}
\]
Quod si autem $b$ ponatur $=\infty$, casus oritur singularis; æquatio namque prodit hæc $dx=-\frac{cdy}{\sqrt{yy-cc}}$, quæ est pro curva Catenaria convexitatem axi $AZ$ obvertente.

\subsection*{Exemplum V.}
\originalsection{45}\sourcepage{196}Inter omnes curvas $az$ æquales areas $aAZz$ continentes, determinare eam, quæ circa axem $AZ$ rotata generet solidum minimæ superficiei.

Quoniam proprietas communis in area $=\int ydx$ constituitur; erit ejus valor differentialis $=nv.dx$. Deinde formula, quæ minimum esse debet, est $\int ydx\sqrt{1+pp}$, cujus valor differentialis est $=nv.\left(dx\sqrt{1+pp}-d.\frac{yp}{\sqrt{1+pp}}\right)$; unde orietur, pro curva quæsita, ista æquatio
\[ndx=dx\sqrt{1+pp}-d.\frac{yp}{\sqrt{1+pp}},\]
quæ, per $p$ multiplicata \& integrata, præbet
\[ny+b=\frac y{\sqrt{1+pp}},\quad\text{seu}\quad\sqrt{1+pp}=\frac y{ny+b};\]
unde fit $p=\frac{\sqrt{y^2-(ny+b)^2}}{ny+b}=\frac{dy}{dx}$; ac
\[dx=\frac{(ny+b)dy}{\sqrt{(1-n^2)y^2-2bny-bb}}.\]
Ex qua patet, si sit $b=0$, tum curvam esse abituram in lineam rectam puncta $a$ \& $z$ jungentem. Deinde si sit $n=0$, ob $dx=\frac{bdy}{\sqrt{yy-bb}}$, curva erit Catenaria concavitatem axi $AZ$ obvertens. Quod si autem sit $n=-1$, fiet $dx=\frac{(b-y)dy}{\sqrt{2by-bb}}$; ex qua integrando oritur $x=c+\frac{2b-y}{3b}\sqrt{2by-bb}$; quæ est pro curva algebraïca, \& in rationalibus præbet $9b(x-c)^2=(2b-y)^2(2y-b)$. Est ideo linea tertii ordinis \& pertinet ad speciem 68 Newtoni.

\subsection*{Exemplum VI.}
\originalsection{46}Inter omnes curvas $az$ ejusdem longitudinis; definire eam quæ circa axem $AZ$ rotata producat maximum solidum.

\sourcepage{197}Inter omnes igitur curvas proprietate communi $\int dx\sqrt{1+pp}$ gaudentes, ea quæritur in qua sit $\int yydx$ maximum. Quoniam ergo formulæ $\int dx\sqrt{1+pp}$ valor differentialis est $=-nv.d.\frac p{\sqrt{1+pp}}$; formulæ vero $\int yydx$ valor differentialis est $=2nv.ydx$; habebitur pro curva quæsita ista æquatio
\[2ydx=\pm bbd.\frac p{\sqrt{1+pp}},\]
quæ, multiplicata per $p$ \& integrata, dabit
\[yy+bc=\pm\frac{bb}{\sqrt{1+pp}};\quad\text{seu}\quad\sqrt{1+pp}=\frac{\pm bb}{yy+bc};\]
hincque $p=\frac{\sqrt{b^4-(yy+bc)^2}}{yy+bc}=\frac{dy}{dx}$; ex qua fit $x=\int\frac{(yy+bc)dy}{\sqrt{b^4-(yy+bc)^2}}$. Hæc curva hanc habet proprietatem ut ejus radius osculi, qui generaliter est $=dx:d.\frac p{\sqrt{1+pp}}$, fiat $=\frac{bb}{2y}$; hoc est, proportionalis est applicatæ $y$ inversæ; unde patet curvam quæsitam esse Elasticam. Non solum autem per constantes $b$ \& $c$ arbitrarias effici potest, ut curva per datos terminos $a$ \& $z$ transeat, sed etiam ut ejus arcus intra hos terminos interceptus fiat datæ magnitudinis. Si sit $c=0$, prodit Elastica rectangula. Cæterum nullo casu constructio per quadraturam vel Circuli vel Hyperbolæ absolvi potest; nisi sint vel $b$ \& $c$ infinita, quo quidem casu linea $az$ prodit recta, vel $b=c$. Hoc enim casu, habebitur
\[x=\int\frac{(yy+bb)dy}{y\sqrt{-(2bb-yy)}},\]
seu, sumto $bb$ negativo, erit
\[
x=\int\frac{(yy-bb)dy}{y\sqrt{2bb-yy}}=f-\sqrt{2bb-yy}-bb\int\frac{dy}{y\sqrt{2bb-yy}};
\]
\& integratione per logarithmos absoluta, fiet
\[x=f-\sqrt{2bb-yy}+bl\frac{b+\sqrt{2bb-yy}}y.\]
Ipsa vero curvæ longitudo, quæ generaliter est $=\int\frac{bbdy}{\sqrt{b^4-(yy+bc)^2}}$, erit hoc casu $=g-bl\frac{b+\sqrt{2bb-yy}}y$.

\subsection*{Exemplum VII.}
\originalsection{47}\sourcepage{198}Invenire curvam, quæ, inter omnes alias ejusdem longitudinis, circa axem $AZ$ rotata, producat solidum cujus superficies sit vel maxima vel minima.

Quoniam proprietas communis est $=\int dx\sqrt{1+pp}$; cujus valor differentialis est $-nv.d.\frac p{\sqrt{1+pp}}$; maximi minimive formulæ autem $\int ydx\sqrt{1+pp}$ valor differentialis est $=nv.\left(dx\sqrt{1+pp}-d.\frac{yp}{\sqrt{1+pp}}\right)$; habebitur pro curva quæsita ista æquatio
\[bd.\frac p{\sqrt{1+pp}}=dx\sqrt{1+pp}-d.\frac{yp}{\sqrt{1+pp}},\]
quæ per $p$ multiplicata \& integrata præbet
\[c-\frac b{\sqrt{1+pp}}=\frac y{\sqrt{1+pp}},\quad\text{seu}\quad c=\frac{b+y}{\sqrt{1+pp}}.\]
Hinc fiet $\sqrt{1+pp}=\frac{b+y}c$, \& $p=\frac{\sqrt{(b+y)^2-cc}}c=\frac{dy}{dx}$: ex hacque $dx=\frac{cdy}{\sqrt{(b+y)^2-cc}}$; quæ est æquatio generalis pro Catenaria, \& satisfacit, dummodo axis respectu catenæ suspensæ situm teneat horizontalem. Fieri igitur potest, ut curva vel convexitatem vel concavitatem axi $AZ$ obvertat, priori casu superficies solidi fiet minima, posteriori maxima.

\subsection*{Exemplum VIII.}
\originalsection{48}\marginnote{\footnotesize \hyperlink{fig:17}{Fig.~17.}}Inter omnes curvas per puncta $A$ \& $C$ transeuntes, quæ omnes æquales areas $ABC$ comprehendant; definire eam quæ in fluido secundum directionem axis $BA$ mota minimam patiatur resistentiam.

Positis abscissa $AP=x$, applicata $PM=y$, proprietas communis est $\int ydx$, ejusque valor differentialis $=nv.dx$. Resistentia autem totalis, quam figura in directione $AB$ sentit, est
\sourcepage{199}ut $\int\frac{p^3dx}{1+pp}$, cujus valor differentialis $-nv.d.\frac{3pp+p^4}{(1+pp)^2}$. Ex his emergit pro curva ista æquatio $dx=bd.\frac{3pp+p^4}{(1+pp)^2}$; quæ integrata dat $x=c+\frac{bpp(3+pp)}{(1+pp)^2}$. Æquatio autem differentialis per $p$ multiplicata, abit in hanc $dy=bpd.\frac{3pp+p^4}{(1+pp)^2}$, quæ in hanc formam
\[
dy=bpd\frac{3pp+p^4}{(1+pp)^2}+bdp.\frac{3pp+p^4}{(1+pp)^2}-bdp.\frac{3pp+p^4}{(1+pp)^4}
\]
transmutata, habet integrale
\[
y=f+\frac{bp^3(3+pp)}{(1+pp)^2}-\frac{bp^3}{1+pp}\quad\text{seu}\quad y=f+\frac{2bp^3}{(1+pp)^2};
\]
cum igitur sit $x=c+\frac{bpp(3+pp)}{(1+pp)^2}$, curva erit algebraïca. Efficiendum est autem, ut, quo casu fit $x=0$ [quod fieri nequit, nisi vel $b$ vel $c$ capiatur negativum] simul $y$ evanescat. Quo autem curva cognoscatur, ponatur $x-c=t$ \& $y-f=u$, erit
\[t=\frac{bpp(3+pp)}{(1+pp)^2},\qquad u=\frac{2bp^3}{(1+pp)^2};\]
unde fit
\[
t+u\sqrt3=\frac{b(p^4+2p^3\sqrt3+3pp)}{(1+pp)^2},\qquad
t-u\sqrt3=\frac{b(p^4-2p^3\sqrt3+3pp)}{(1+pp)^2}.
\]
Extrahendis igitur radicibus quadratis habebitur
\[
\sqrt{\frac{t+u\sqrt3}{b}}=\frac{pp+p\sqrt3}{b},\qquad
\sqrt{\frac{t-u\sqrt3}{b}}=\frac{pp-p\sqrt3}{1+pp};
\]
hincque
\[
\begin{aligned}
\sqrt{\frac{t+u\sqrt3}{b}}+\sqrt{\frac{t-u\sqrt3}{b}}&=\frac{2pp}{1+pp},\\
\sqrt{\frac{t+u\sqrt3}{b}}-\sqrt{\frac{t-u\sqrt3}{b}}&=\frac{2p\sqrt3}{1+pp}.
\end{aligned}
\]
At est
\[
\begin{aligned}
\frac tb&=\frac32.\frac{2pp}{1+pp}-\frac12.\frac{4p^4}{(1+pp)^2}\\
&=\frac32\sqrt{\frac{t+u\sqrt3}{b}}+\frac32\sqrt{\frac{t-u\sqrt3}{b}}
-\frac12\left(\frac{2t}b+2\sqrt{\frac{tt-3uu}{bb}}\right).
\end{aligned}
\]
Ergo
\[\frac{4t}b=3\sqrt{\frac{t+u\sqrt3}{b}}+3\sqrt{\frac{t-u\sqrt3}{b}}-\frac{2\sqrt{tt-3uu}}b;\]
quæ rationalis facta præbet æquationem hanc quarti ordinis
\sourcepage{200}
\[
4t^4+8ttuu+4u^4=4bt^3+36btuu-27b^2uu,
\]
seu $4(tt+uu)^2=4bt^3+36btu^2-27b^2u^2$.

Ad curvam autem per infinita puncta construendam expedit adhibere has formulas, $t=\frac{b(p^4+3pp)}{(1+pp)^2}$ \& $u=\frac{2bp^3}{(1+pp)^2}$. Primum autem patet curvam habere diametrum in positione abscissarum $t$ sitam, duobusque locis fieri $u=0$, nempe casu $p=0$, quo simul fit $t=0$, \& casu $p=\infty$, quo fit $t=b$. Quod si ponatur $b=4c$, atque $t=3c+r$, orietur ista æquatio
\[(rr+uu)^2+8c(r^3-3ru^2)+18cc(r^2+u^2)-27c^4=0\]
quæ cum sit functio ipsarum $rr+uu$ \& $r^3-3ruu$, declarat curvam hanc habere tres diametros sese in initio abscissarum harum $r$ decussantes. Curva ergo quæsita triangulo æquilatero $ABC$ ita erit inscriptibilis, ut constet ex tribus ramis $ADB$, $AEC$ \& $BFC$ inter se similibus \& æqualibus, qui in punctis $A$, $B$, \& $C$ cuspides forment acutissimos. Ejus igitur diametri erunt tres rectæ $AI$, $BH$ \& $CG$, sese in centro trianguli $O$ decussantibus. Erit autem $AO=3c$, $OF=c$, \& $OI=\frac32c$, ita ut sit $AI=\frac92c$ \& $FI=\frac12c=\frac12OF$. Hujus jam curvæ quæcunque portio $abc$ rectis $ab$ \& $bc$ parallelis ipsis $AI$ \& $BI$ \& arcu curvæ $ac$ comprehensa, ita erit comparata, ut arcus $ac$ inter omnes alios puncta $a$ \& $c$ jungentes, \& æqualem aream $abc$ continentes, in fluido secundum directionem $ba$ mota minimam patiatur resistentiam. Porro autem hæc curva erit rectificabilis, reperiturque arcus $ADB=\frac{16}3c$; ex quo erit $ADB:AI=\frac{16}3:\frac92=32:27$, atque $ADB:AB=32:18\sqrt3=16:9\sqrt3$.

\subsection*{Exemplum IX.}
\originalsection{49}\marginnote{\footnotesize \hyperlink{fig:19}{Fig.~19.}}Inter omnes curvas $AM$ æquales areas $APM$ includentes; invenire eam, quæ sit ita comparata, ut si perpetuo a centro circuli osculantis $O$ ad applicatam $MP$ productam ducatur perpendicularis $ON$; curva a punctis $N$ formata minimam comprehendat aream $APN$.

Positis abscissa $AP=x$, \& applicata $PM=y$; erit area
\sourcepage{201}$APM=\int ydx$, quæ est proprietas communis, ejusque valor differentialis $=nv.dx$. Deinde, cum sit radius osculi $MO=-\frac{(1+pp)^{3:2}}q$, fiet $MN=-\frac{1+pp}q$ \& $PN=-\frac{1+pp}q-y$; ex quo area $APN$ erit $=-\int ydx-\int\frac{1+pp}qdx$; quæ debet esse minima, cujus valor differentialis est
\[=nv.\left(-dx+d.\frac{2p}q+\frac1{dx}dd.\frac{1+pp}{qq}\right);\]
unde ista nascitur æquatio
\[ndx^2=dxd.\frac{2p}q+dd.\frac{1+pp}{qq};\]
quæ integrata dat $nxdx=\frac{2pdx}q+d.\frac{1+pp}{qq}+bdx$. Illa vero eadem æquatio, per $p$ multiplicata, dat $ndxdy=dyd.\frac{2p}q+pdd.\frac{1+pp}{qq}$; cujus integrale est $nydx=cdx-\frac{2dx}q+pd.\frac{1+pp}{qq}$. His æquationibus conjungendis, oritur
\[
\begin{aligned}
nxdy-nydx&=bdy-cdx+\frac{2pdy}q+\frac{2dx}q\\
&=bdy-cdx+\frac{2dx^2+2dy^2}{dp}.
\end{aligned}
\]
Ponatur $nx-b=nt$, \& $ny-c=nu$; erit $dy=du$, \& $dx=dt$, atque
\[ndp=\frac{2dt^2+2du^2}{tdu-udt}=\frac{nddu}{dt},\]
seu $2dt^3+2dtdu^2=ntduddu-nudtddu$, posito $dt$ constante. Sit $u=st$, erit $du=sdt+tds$, \& $ddu=tdds+2dtds$; hisque substitutis prodibit ista æquatio:
\[2(1+ss)dt^3+4stdt^2ds+2(1-n)ttdtds^2=nt^3dsdds.\]
Ponatur $t=e^{\int rds}$, erit $dt=e^{\int rds}rds$, \& $ddt=0=e^{\int rds}(rdds+drds+rrds^2)$; unde fit $dds=-\frac{drds}r-rds^2$; ex quibus tandem emergit
\[2(1+ss)r^3ds+4sr^2ds+2(1-n)rds=-\frac{ndr}r-nrds,\]
seu
\sourcepage{202}
\[\frac{ndr}r+(2-n)rds+4sr^2ds+2r^3ds+2r^3s^2ds=0.\]
Sit $s=v-\frac1r$, fiet $dr+rrdv=\frac{ndv}{2(1+vv)}$; quæ æquatio integrationem admittit, quoties est $n=2i(i-1)$ denotante $i$ numerum integrum quemcunque: ut si sit $n=4$, fiet
\[r=\frac{2v}{1+vv}+\frac1{(1+vv)^2\int\frac{dv}{(1+vv)^2}};\]
ex qua retrogradiendo constructio absolvi poterit.

\subsection*{Exemplum X.}
\originalsection{50}Inter omnes curvas, in quibus $\int xTdx$ eundem obtinet valorem; invenire eam in qua sit $\int yTdx$ maximum vel minimum, existente $T$ functione quacunque ipsius $p$, ita ut sit $dT=Pdp$.

Ad formulæ $\int xTdx$ valorem differentialem inveniendum, notandum est esse $d.xT=Tdx+xPdp$, ex quo illius valor differentialis erit $=-nv.d.xP$. Ex altera autem formula $\int yTdx$ habetur $d.yT=Tdy+yPdp$, unde ejus valor differentialis erit $nv.(Tdx-d.yP)$. Quare, pro curva quæsita orietur ista æquatio $nd.xP=Tdx-d.yP$. Ergo $\int Tdx=nxP+yP+b$. Porro si illa æquatio per $p$ multiplicetur, habebitur
\[npd.xP=Tdy-pd.yP=d.yT-yP.dp-pd.yP=d.yT-d.yPp.\]
At est
\[pd.xP=Ppdx+pxdP+xPdp+Tdx-d.xT=d.xPp+Tdx-dxT.\]
Quamobrem orietur $d.yT-d.yPp=nd.xPp+nTdx-nd.xT$; hincque $\int nTdx=yT-yPp-nxPp+nxT+c$. Quia vero, ex superiori integratione habemus $\int nTdx=nnxP+nyP+nb$, erit, eliminando $\int nTdx$, ista æquatio $nnxP+nyP+nb=yT-yPp-nxPp+nxT+c$, seu
\[y=\frac{nx(nP+Pp-T)+c}{-nP-Pp+T},\quad\text{vel}\quad y=-nx+\frac c{T-nP-Pp}.\]
\sourcepage{203}Ergo prodit tandem
\[
\begin{aligned}
x&=c\int\frac{dP}{(T-nP-Pp)^2},\\
y&=c\int\frac{pdP}{(T-nP-Pp)^2}=\frac c{T-nP-Pp}-nc\int\frac{dP}{(T-nP-Pp)^2}.
\end{aligned}
\]

\subsection*{Exemplum XI.}
\originalsection{51}Invenire curvam, quæ, inter omnes alias intra eosdem terminos contentas, \& eundem formulæ $\int xdx\sqrt{1+pp}$ valorem continentes, habeat $\int ydx\sqrt{1+pp}$ maximum vel minimum.

Exemplum hoc est Casus præcedentis, atque ex illo manat, ponendo $T=\sqrt{1+pp}$, ex quo erit $P=\frac p{\sqrt{1+pp}}$, \& $dP=\frac{dp}{(1+pp)^{3:2}}$. Porro vero erit $T-nP-Pp=\frac{1-np}{\sqrt{1+pp}}$. Ex his jam surrogatis, prodibit
\[x=c\int\frac{dp}{(1-np)^2\sqrt{1+pp}},\qquad y=\frac{c\sqrt{1+pp}}{1-np}-nx.\]
Integratione autem per logarithmos instituta fiet
\[
\begin{aligned}
x={}&\frac{nc(p+\sqrt{1+pp})-c}{(1+nn)(1-np)}\\
&+\frac c{(1+nn)^{3:2}}l\frac{n+(1+\sqrt{1+nn})(p+\sqrt{1+pp})}{n+(1-\sqrt{1+nn})(p+\sqrt{1+pp})}+b,
\end{aligned}
\]
\&
\[
\begin{aligned}
y={}&\frac{nc+c(\sqrt{1+pp}-nnp)}{(1+nn)(1-np)}\\
&-\frac{nc}{(1+nn)^{3:2}}l\frac{n+(1+\sqrt{1+nn})(p+\sqrt{1+pp})}{n+(1-\sqrt{1+nn})(p+\sqrt{1+pp})}-nb;
\end{aligned}
\]
ex quibus valoribus curva construi poterit per logarithmos. Generaliter autem, quamcunque $T$ functionem ipsius $p$ denotet, constructio semper per quadraturas absolvi potest. Cæterum hoc Exemplum sine subsidio præcedentis multo difficilius solutu fuisset\sourcepage{204}; non tam facile enim perspicere licuisset, quomodo æquatio inventa integrabilis redderetur quam in casu generali.

\subsection*{Propositio V. Problema.}
\originalsection{52}Inter omnes curvas ad eandem abscissam $=a$ relatas, quæ eundem formulæ $\pi=\int[Z]dx$ valorem recipiunt; invenire eam, in qua sit $\int Zdx$ maximum vel minimum; existente $Z$ functione simul ipsius $\pi$, ita ut sit $dZ=Ld\pi+Mdx+Ndy+Pdp+Qdq+\&c.$ atque $d[Z]=[M]dx+[N]dy+[P]dp+[Q]dq+\&c.$

\subsection*{Solutio.}
Quoniam est $d[Z]=[M]dx+[N]dy+[P]dp+[Q]dq+\&c.$ erit formulæ $\int[Z]dx$, quæ hic quantitatem omnibus curvis communem repræsentat, valor differentialis
\[=nv.dx\left([N]-\frac{d[P]}{dx}+\frac{dd[Q]}{dx^2}-\&c.\right),\]
qui ex Casu primo §.~7 Cap. præced. sequitur. At formula $\int Zdx$, maximum minimumve exprimens, quia $Z$ involvit formulam integralem $\pi=\int[Z]dx$, pertinet ad Casum secundum loci citati: ejusque adeo valor differentialis erit
\[=nv.dx\left(N+[N]V-\frac{d(P+[P]V)}{dx}+\frac{dd(Q+[Q]V)}{dx^2}-\&c.\right),\]
denotante $V=H-\int Ldx$, ubi $H$ est quantitas determinata, quæ oritur si in integrali $\int Ldx$ ponatur $x=a$. Atque ob hanc ipsam quantitatem $H$, iste valor differentialis a præscripta longitudine abscissæ $x=a$ pendet. Ex his igitur duobus valoribus differentialibus ambarum formularum propositarum, quarum altera proprietatem communem, altera maximum minimumve exponit, secundum regulam datam, nascitur æquatio pro curva sequens:
\[
\begin{aligned}
0={}&\alpha[N]-\frac{\alpha d[P]}{dx}+\frac{\alpha dd[Q]}{dx^2}-\&c.+N+[N]V\\
&-\frac{d(P+[P]V)}{dx}+\frac{dd(Q+[Q]V)}{dx^2}-\&c.
\end{aligned}
\]
quæ, ob $V=\sourcepage{205}H-\int Ldx$, transit in hanc
\[
\begin{aligned}
0={}&N+(\alpha+H-\int Ldx)[N]\\
&-\frac{d(P+(\alpha+H-\int Ldx)[P])}{dx}\\
&+\frac{dd(Q+(\alpha+H-\int Ldx)[Q])}{dx^2}-\&c.
\end{aligned}
\]
Cum jam $\alpha$ sit quantitas constans arbitraria; etiamsi $H$ sit quantitas constans determinata, tamen $\alpha+H$ fiet quantitas arbitraria: ideoque non amplius a definita abscissæ longitudine $a$ pendet. Quare si, loco $\alpha+H$, scribamus $C$, habebimus pro curva quæsita hanc æquationem:
\[
\begin{aligned}
0={}&N+(C-\int Ldx)[N]-\frac{d(P+(C-\int Ldx)[P])}{dx}\\
&+\frac{dd(Q+(C-\int Ldx)[Q])}{dx^2}-\&c.
\end{aligned}
\]
quæ ergo pro quacunque abscissa exhibet Curvam, quæ, inter omnes alias eundem formulæ $\int[Z]dx$ valorem recipientes, continebit formulæ $\int Zdx$ maximum minimumve valorem. Q. E. I.

\subsection*{Corollarium I.}
\originalsection{53}Si igitur proprietas communis fuerit ea ipsa formula integralis, quæ in maximi minimive formula implicatur; tum consideratio determinatæ abscissæ magnitudinis ex calculo egreditur, \& Curva inventa pro quavis abscissa quæsito satisfacit.

\subsection*{Corollarium II.}
\originalsection{54}In hac æquatione inventa, duæ adhuc inerunt formulæ integrales; primo nempe formula $\int Ldx$, ac deinde formula $\pi=\int[Z]dx$, quæ cum ea in $Z$ contineatur, inerit in quantitatibus $L$, $M$, $N$, $P$ \&c.

\subsection*{Corollarium III.}
\originalsection{55}Si igitur hæc integralia per differentiationem tollere lubeat; pervenietur ad differentialia binis gradibus altiora, simulque exibit constans arbitraria $C$. Interim tamen numerus cons\sourcepage{206}tantium arbitrariarum unitate minor erit quam gradus iste differentialium; eo quod integrale $\pi=\int[Z]dx$ definitum obtinere debet valorem, eum ipsum scilicet, quem in maximi minimive formula $\int Zdx$ habet.

\subsection*{Corollarium IV.}
\originalsection{56}Hinc igitur in æquatione inventa, ob constantem arbitrariam $C$, potestate una plures inerunt constantes, quam differentialium gradus indicat. Quarum una eo determinabitur, ut valor formulæ communis $\pi=\int[Z]dx$ fiat pro curva inventa datæ magnitudinis; reliquæ vero per data puncta vel tangentium positionem datam determinabuntur.

\subsection*{Corollarium V.}
\originalsection{57}Si $Z$ fuerit functio cum quantitatum $x$, $y$, $p$, $q$, \&c. tum arcus curvæ $s$: atque inter omnes Curvas isoperimetras quæratur ea, in qua sit $\int Zdx$ maximum vel minimum; tum fiet $\pi=s=\int[Z]dx$ \& $[Z]=\sqrt{1+pp}$, ita ut sit $[M]=0$, $[N]=0$, \& $[P]=\frac p{\sqrt{1+pp}}$.

\subsection*{Corollarium VI.}
\originalsection{58}Hoc igitur casu, si fuerit $dZ=Lds+Mdx+Ndy+Pdp+Qdq+\&c.$ habebitur pro Curva quæ, inter omnes isoperimetras, habeat $\int Zdx$ maximum vel minimum, ista æquatio:
\[0=N-\frac1{dx}d\left(P+\frac{(C-\int Ldx)p}{\sqrt{1+pp}}\right)+\frac{ddQ}{dx^2}-\&c.\]
seu
\[N-\frac{dP}{dx}+\frac{ddQ}{dx^2}=\frac{(C-\int Ldx)dp}{dx(1+pp)^{3:2}}-\frac{Lp}{\sqrt{1+pp}},\]
sive
\[\frac{Lp}{\sqrt{1+pp}}+N-\frac{dP}{dx}+\frac{ddQ}{dx^2}-\&c.=\frac{(C-\int Ldx)dp}{dx(1+pp)^{3:2}}.\]

\subsection*{Corollarium VII.}
\originalsection{59}\sourcepage{207}Cum sit $C$ quantitas arbitraria; in genere notari convenit, quod si pro $C$ accipiatur ille formulæ $\int Ldx$ valor, quem induit si ponatur $x=a$, tum prodituram esse curvam, quæ inter omnes omnino curvas eidem abscissæ $x=a$ respondentes, habeat valorem formulæ $\int Zdx$ maximum vel minimum.

\subsection*{Scholion I.}
\originalsection{60}Casus Coroll.~6, quia is ab Auctoribus potissimum tractari est solitus, peculiarem evolutionem meretur, ut ejus ope Problemata quæ forte occurrere queunt, facilius \& expeditius resolvi possint. Inter omnes igitur Curvas isoperimetras, seu quæ eandem habeant longitudinem $s=\int dx\sqrt{1+pp}$, quæratur ea, in qua sit $\int Zdx$ maximum vel minimum, existente $Z$ functione cum quantitatum definitarum $x$, $y$, $p$, $q$ \&c. tum arcus curvæ $s$; ita ut sit $dZ=Lds+Mdx+Ndp+Pdp+\&c.$ Pro curva hac proprietate gaudente jam inventa est hæc æquatio:
\[\frac1{dx}d.\frac{(C-\int Ldx)p}{\sqrt{1+pp}}=N-\frac{dP}{dx}+\frac{ddQ}{dx^2}-\&c.\]
quæ quidem, in hoc latissimo sensu nec integrari nec ad simpliciorem formam se reduci patitur. At casus notasse juvabit, quibus eam integrare licebit. Ac primo quidem si sit $N=0$ sponte prodit ista pro curva æquatio:
\[A+\frac{(C-\int Ldx)p}{\sqrt{1+pp}}=-P+\frac{dQ}{dx}-\&c.\]
jam semel integrata. Secundo ponamus esse $M=0$; atque æquatio per $pdx=dy$ multiplicata abibit in hanc
\[pd.\frac{(C-\int Ldx)p}{\sqrt{1+pp}}=Ndy-Pdp+\frac{pddQ}{dx}-\&c.\]
ad quam si addatur $Lds=Ldx\sqrt{1+pp}=dZ-Ndy-Pdp-Qdq$ \&c.; integratione instituta prodibit
\[\int\left(Ldx\sqrt{1+pp}+\frac{pd.(C-\int Ldx)p}{\sqrt{1+pp}}\right)=Z-Pp-Qq+\frac{pdQ}{dx}\quad\&c.\]
\sourcepage{208}Prius vero membrum si evolvatur, transit in
\[
\begin{aligned}
&\int\left(Ldx\sqrt{1+pp}+\frac{(C-\int Ldx)pdp}{(1+pp)^{3:2}}-\frac{Lppdx}{\sqrt{1+pp}}\right)\\
&=\int\left(\frac{Ldx}{\sqrt{1+pp}}+\frac{(C-\int Ldx)pdp}{(1+pp)^{3:2}}\right),
\end{aligned}
\]
cujus integrale est $-\frac{C-\int Ldx}{\sqrt{1+pp}}$. Quare, casu quo $M=0$, habebitur ista æquatio
\[\frac{C-\int Ldx}{\sqrt{1+pp}}=A-Z+Pp+Qq-\frac{pdQ}{dx}.\]
Sin autem tertio fuerit tam $M=0$ quam $N=0$, habebitur primum, ob $N=0$, hæc æquatio:
\[A+\frac{(C-\int Ldx)p}{\sqrt{1+pp}}=-P+\frac{dQ}{dx};\]
quæ, multiplicata per $dp=qdx$, abit in hanc
\[Adp+\frac{(C-\int Ldx)pdp}{\sqrt{1+pp}}=-Pdp+Qdq.\]
Cum autem sit $dZ=Ldx\sqrt{1+pp}+Pdp+Qdq$, habebitur
\[dZ+Adp-Ldx\sqrt{1+pp}+\frac{(C-\int Ldx)pdp}{\sqrt{1+pp}}=qdQ+Qdq;\]
quæ integrata dabit, $Z+B+Ap+(C-\int Ldx)\sqrt{1+pp}=Qq$, seu
\[C-\int Ldx=\frac{Qq-B-Ap-Z}{\sqrt{1+pp}}.\]
At ex priore æquatione est
\[C-\int Ldx=-\frac{A\sqrt{1+pp}}p-\frac{P\sqrt{1+pp}}p+\frac{dQ\sqrt{1+pp}}{pdx};\]
ex quibus conjungendis elicitur:
\[Adx-Bdy=Zdy-Pdx-Ppdy+dQ+ppdQ-Qpdp,\]
in qua non amplius inest formula integralis $\int Ldx$. Usum igitur horum casuum in Exemplis monstrabimus.

\subsection*{Exemplum I.}
\originalsection{61}Inter omnes curvas isoperimetras; definire eam, in qua sit $\int s^n dx$ maximum vel minimum, denotante $s$ arcum curvæ abscissæ $x$ respondentem.

\sourcepage{209}Quoniam proprietas communis longitudinem arcus $s=\int dx\sqrt{1+pp}$ respicit, atque in maximi minimive formula $\int s^n dx$ inest ipse arcus, solutio pertinebit ad Casum in Scholio pertractatum. Comparata ergo formula $\int s^n dx$ cum generali $\int Zdx$, fiet $Z=s^n$ \& $dZ=ns^{n-1}ds$ hincque $L=ns^{n-1}$, $M=0$, $N=0$, $P=0$ \&c. Quare ex Scholii Casu ultimo, quo posueramus $M=0$ \& $N=0$, habebitur ista æquatio $Adx-Bdy=Zdy=s^n dy$, ex qua oritur
\[
Adx=dy(B+s^n),\qquad A^2dx^2+A^2dy^2=A^2ds^2=dy^2(A^2+(B+s^n)^2)
\]
ideoque
\[dy=\frac{Ads}{\sqrt{A^2+(B+s^n)^2}},\qquad dx=\frac{(B+s^n)ds}{\sqrt{A^2+(B+s^n)^2}};\]
unde Curvæ constructio perfici poterit. Vel posito $dy=pdx$, erit $s^n=\frac{A-Bp}p$; atque $s=\sqrt[n]{\frac{A-Bp}p}$: ex quo fiet
\[ds=dx\sqrt{1+pp}=-\frac{Adp(A-Bp)^{(1-n):n}}{np^{(1+n):n}}.\]
Atque hinc per $p$ coordinatæ curvæ $x$ \& $y$ ita determinabuntur, ut sit
\[
x=-\frac An\int\frac{dp(A-Bp)^{(1-n):n}}{p^{(1+n):n}\sqrt{1+pp}},\qquad
y=-\frac An\int\frac{dp(A-Bp)^{(1-n):n}}{p^{1:n}\sqrt{1+pp}}.
\]
Videntur hic quidem quatuor constantes, duæ scilicet novæ, præter $A$ \& $B$, ingredi, ob duplicem integrationem $y$ \& $x$. At cum posito $x=0$, simul arcus curvæ $s=\sqrt[n]{\frac{A-Bp}p}$ evanescere debeat; hinc vicissim constans in integratione ipsius $x$ orta definitur. Nimirum si $n$ fuerit numerus affirmativus, arcus\sourcepage{210} $s$ evanescit, posito $p=\frac AB$; ex quo valor ipsius $x$ ita determinari debet, ut posito $p=\frac AB$ fiat $=0$.

Quod si ponatur $n=1$; habebitur ex priore constructione, statim $dx=\frac{(B+s)ds}{\sqrt{A^2+(B+s)^2}}$; ideoque $x=\sqrt{A^2+B^2+2Bs+ss}-\sqrt{A^2+B^2}$, seu posito $B=b$, \& $\sqrt{A^2+B^2}=c$, erit $x+c=\sqrt{c^2+2bs+ss}$. Ex posteriore autem construendi modo, oritur
\[x=-A\int\frac{dp}{pp\sqrt{1+pp}}=\frac{A\sqrt{1+pp}}p+b,\]
seu $(x-b)p=c\sqrt{1+pp}$; hincque $p=\frac c{\sqrt{(x-b)^2-c^2}}=\frac{dy}{dx}$. Quare cum sit $y=\int\frac{cdx}{\sqrt{(x-b)^2-c^2}}$; curva satisfaciens erit Catenaria.

\subsection*{Exemplum II.}
\originalsection{62}Inter omnes curvas ejusdem longitudinis, eam determinare, in qua sit $\int Sdx$ maximum vel minimum, existente $S$ functione quacunque arcus $s$.

Quia proprietas communis arcu $s=\int dx\sqrt{1+pp}$ continetur; solutio ex Scholio peti poterit. Scilicet cum sit $Z=S=$ functioni ipsius $s$, erit $Lds=dS$, \& $M=N=P=Q$ \&c. $=0$. Quare, per tertium Scholii Casum, habebitur pro curva quæsita ista æquatio $Adx-Bdy=Sdy$, \& $Adx=dy(B+S)$. Hinc ergo erit
\[A^2dx^2+A^2dy^2=A^2ds^2=dy^2(A^2+(B+S)^2)\]
\& $y=\int\frac{Ads}{\sqrt{A^2+(B+S)^2}}$; erit autem abscissa $x=\int\frac{(B+S)ds}{\sqrt{A^2+(B+S)^2}}$; unde curvæ constructio absolvi poterit.

Ponamus esse $S=e^s$; positoque $dy=pdx$, erit $\frac{A-Bp}p\sourcepage{211}=e^s$, \&
\[e^sds=-\frac{Adp}{pp}=\frac{(A-Bp)dx\sqrt{1+pp}}p,\]
hincque
\[dx=\frac{-Adp}{(A-Bp)p\sqrt{1+pp}},\qquad dy=\frac{-Adp}{(A-Bp)\sqrt{1+pp}}.\]
Componendo vero fiet $dx-\frac{Bdy}A=\frac{-dp}{p\sqrt{1+pp}}$, \& integrando $Ax-By=Al\frac{1+\sqrt{1+pp}}p+C$, seu $\frac{1+\sqrt{1+pp}}p=e^{(Ax-By-C):A}$. Cum autem, facto $s=0$, evanescere debeat $x$, atque ob $\frac{A-Bp}p=e^s$, facto $s=0$, fiat $p=\frac A{A+1}$, per integrationes efficiendum est, ut, facto $p=\frac A{B+1}$ fiat $x=0$.

\subsection*{Exemplum III.}
\originalsection{63}Inter omnes curvas ejusdem longitudinis, determinare eam in qua sit $\int sydx$ maximum vel minimum, denotante $s$ arcum curvæ.

Solutio hujus Quæstionis iterum petenda est ex Scholio; erit namque $Z=sy$ \& $dZ=yds+sdy$, ex quo fit $L=y$, $M=0$ \& $N=s$, reliquæ litteræ $P$, $Q$, \&c. evanescent. Cum igitur sit $M=0$, Casus Scholii secundus hanc suppeditabit solutionem: $\frac{C-\int ydx}{\sqrt{1+pp}}=A-ys$; immediate vero prodit
\[sdx=d.\frac{(C-\int ydx)p}{\sqrt{1+pp}}=\frac{(C-\int ydx)dp}{(1+pp)^{3:2}}-\frac{ypdx}{\sqrt{1+pp}}.\]
Quare, cum sit $C-\int ydx=A\sqrt{1+pp}-ys\sqrt{1+pp}$, erit $sdx=\frac{Adp-ysdp-ydy}{1+pp}$, seu $sdx+spdy+ysdp+ydy=Adp$. Sin autem libuerit arcum $s$ eliminare, habebitur ex binis æquationibus,
\[s=\frac Ay-\frac{C-\int ydx}{y\sqrt{1+pp}}=\frac{(C-\int ydx)dp}{dx(1+pp)^{3:2}}-\frac{yp}{\sqrt{1+pp}};\]
\sourcepage{212}hincque
\[\frac{Adx}y+\frac{ypdx}{\sqrt{1+pp}}=(C-\int ydx)\left(\frac{dp}{(1+pp)^{3:2}}+\frac{dx}{y\sqrt{1+pp}}\right).\]
In utroque autem casu difficile est ad æquationem ad curvam construendum accommodatam pertingere.

\subsection*{Exemplum IV.}
\originalsection{64}Inter omnes curvas eandem aream $\pi=\int ydx$ continentes, definire eam, in qua sit $\int\frac{dx\sqrt{1+pp}}\pi$ maximum vel minimum.

Si hanc Quæstionem cum Solutione generali comparemus, habebimus $\int[Z]dx=\int ydx$; hincque $[Z]=y$, \& $[N]=1$; reliquis litteris $[M]$, $[P]$, $[Q]$, \&c. evanescentibus. Porro erit
\[Z=\frac{\sqrt{1+pp}}\pi,\qquad dZ=-\frac{d\pi\sqrt{1+pp}}{\pi^2}+\frac{pdp}{\pi\sqrt{1+pp}},\]
unde erit $L=-\frac{\sqrt{1+pp}}{\pi^2}$, $M=0$, $N=0$, \& $P=\frac p{\pi\sqrt{1+pp}}$. Quocirca pro curva quæsita sequens emerget æquatio:
\[0=C+\int\frac{dx\sqrt{1+pp}}{\pi^2}-\frac1{dx}d.\frac p{\pi\sqrt{1+pp}}.\]
Multiplicetur hæc æquatio per $dy=pdx$, erit
\[0=Cdy+dy\int\frac{dx\sqrt{1+pp}}{\pi^2}-pd.\frac p{\pi\sqrt{1+pp}},\]
quæ integrata dabit:
\[
\begin{aligned}
0={}&B+Cy+y\int\frac{dx\sqrt{1+pp}}{\pi^2}-\int\frac{d\pi\sqrt{1+pp}}{\pi^2}\\
&-\frac{pp}{\pi\sqrt{1+pp}}+\int\frac{pdp}{\pi\sqrt{1+pp}}\\
={}&B+Cy+y\int\frac{dx\sqrt{1+pp}}{\pi^2}-\frac{pp}{\pi\sqrt{1+pp}}+\frac{\sqrt{1+pp}}\pi.
\end{aligned}
\]
Hinc itaque istam obtinebimus æquationem
\[0=B+Cy+y\int\frac{dx\sqrt{1+pp}}{\pi^2}+\frac1{\pi\sqrt{1+pp}};\]
a qua si prior per $y$ multiplicata subtrahatur, erit
\sourcepage{213}
\[0=B+\frac1{\pi\sqrt{1+pp}}+\frac y{dx}d.\frac p{\pi\sqrt{1+pp}},\]
seu
\[0=Bdx+\frac{dx}{\pi\sqrt{1+pp}}+\frac{ydp}{\pi(1+pp)^{3:2}}-\frac{y^2pdx}{\pi^2\sqrt{1+pp}};\]
ex qua æquatione si denuo $\pi=\int ydx$ exterminare velimus, prodiret æquatio differentialis tertii ordinis, ex qua multo minus quicquam ad Curvam cognoscendam deduci posset.

\subsection*{Scholion II.}
\originalsection{65}Quanquam, in hac Propositione posuimus $[Z]$ esse functionem determinatam quantitatum $x$, $y$, $p$, $q$, \&c. tamen Methodus solvendi patet, si hæc ipsa quantitas $[Z]$ fuerit functio indefinita formulas integrales in se complectens. Ponamus enim in formula $\pi=\int[Z]dx$, quæ omnibus curvis debet esse communis, esse
\[d[Z]=[L]d\varpi+[M]dx+[N]dy+[P]dp+[Q]dq+\&c.\]
existente $\varpi=\int[z]dx$, \&
\[d[z]=[m]dx+[n]dy+[p]dp+[q]dq+\&c.\]
Maximum minimumve autem esse oportere formulam $\int Zdx$, existente: $dZ=Ld\pi+Mdx+Ndy+Pdp+\&c.$ Jam formula $\int[Z]dx$ continetur in Casu secundo §.~7 Cap. præc: inde ergo si capiatur integrale $\int[L]dx$ ejusque valor respondens abscissæ $x=a$, ad quam solutio debet accommodari, ponatur $=[H]$, atque $[H]-\int[L]dx=[V]$; habebitur formulæ $\int[Z]dx$ valor differentialis
\[=nv.dx\left([N]+[n][V]-\frac{d([P]+[p][V])}{dx}+\frac{dd([Q]+[q][V])}{dx^2}-\&c.\right).\]
Deinde vero maximi minimive formula $\int Zdx$ continetur in Casu tertio loci citati; ad ejusque valorem differentialem inveniendum, ponatur formulæ $\int Ldx$ valor abscissæ $x=a$ respondens $=H$, ac $H-\int Ldx=V$. Jam capiatur integrale
\[\int[L]Vdx=H\int[L]dx-\int[L]dx\int Ldx\]
sitque, posito $x=a$, valor formulæ $\int[L]dx\int Ldx=K$, eodem autem casu formulæ $\int[L]dx$ valor est $=[H]$, ex quo formulæ
\sourcepage{214}$\int[L]Vdx$, casu $x=a$, valor erit $=H[H]-K$, \& vocetur
\[H[H]-K-H\int[L]dx+\int[L]dx\int Ldx=W\]
ita ut sit $W=H[V]-K+\int[L]dx\int Ldx$, eritque formulæ propositæ $\int Zdx$ valor differentialis
\[
=nv.dx\left(N+[N]V+[n]W-\frac{d(P+[P]V+[p]W)}{dx}+\frac{dd(Q+[Q]V+[q]W)}{dx^2}-\&c.\right).
\]
Quod si jam ad hunc valorem differentialem addatur præcedens per quantitatem constantem arbitrariam $\alpha$ multiplicatus, summaque ponatur $=0$, prodibit æquatio pro curva quæsita hæc:
\[
\begin{aligned}
0={}&N+[N](\alpha+V)+[n](\alpha[V]+W)\\
&-\frac1{dx}d(P+[P](\alpha+V)+[p](\alpha[V]+W))\\
&+\frac1{dx^2}dd(Q+[Q](\alpha+V)+[q](\alpha[V]+W))-\&c.
\end{aligned}
\]
Est vero hic $\alpha+V=\alpha+H-\int Ldx$; unde si ponatur $\alpha+H=C$, erit $C$ constans arbitraria, \& $\alpha+V=C-\int Ldx$, atque
\[\alpha[V]+W=C[H]-K-C\int[L]dx+\int[L]dx\int Ldx.\]
Hoc igitur pacto, pervenietur ad curvam quæsitam, in cujus æquatione, quia ob $[H]$ \& $K$ adhuc inest constans data $a$, ea quæsito satisfaciet tantum pro proposita abscissa $x=a$. Quod si autem formularum ambarum altera ad Casum 4, altera ad Casum 5 pertineat, tum iterum consideratio datæ abscissæ $a$ ex calculo egreditur, eademque curva pro omni abscissa satisfaciet, id quod unico sequenti Exemplo declarasse sufficiet.

\subsection*{Exemplum V.}
\originalsection{66}Inter omnes curvas eidem abscissæ respondentes, quæ eundem formulæ $v$ valorem recipiunt; invenire eam, in qua sit $\int\frac{dx\sqrt{1+pp}}{\sqrt v}$ maximum vel minimum, existente $dv=gdx+Wdx\sqrt{1+pp}$ \& $W$ functione quacunque ipsius $v$.

Solutio hujus Quæstionis exhibebit curvam super qua corpus
\sourcepage{215}descendens a gravitate uniformi $g$ deorsum, in directione abscissarum sollicitatum in medio quocunque resistente celerrime delabitur, inter omnes alias curvas super quibus descendendo eandem acquirit celeritatem. Est enim $\sqrt v$ celeritas corporis in quocunque curvæ puncto, \& $W$ exprimit resistentiam medii. Quod nunc primum ad proprietatem communem $v=\int dx(g+W\sqrt{1+pp})$; ponamus esse $dW=Udv$, atque hæc formula ad Casum quartum pertinebit; erit namque $\pi=v$, \& $Z=g+W\sqrt{1+pp}$, ac
\[dZ=Udv\sqrt{1+pp}+\frac{Wpdp}{\sqrt{1+pp}};\]
unde erit $L=U\sqrt{1+pp}$, $M=0$, $N=0$, \& $P=\frac{Wp}{\sqrt{1+pp}}$. Sumatur ergo integrale $\int Udx\sqrt{1+pp}$ sitque, casu quo $x=a$ ponitur, $e^{\int Udx\sqrt{1+pp}}=H$, ac ponatur $V=He^{-\int Udx\sqrt{1+pp}}$. Ex his erit formulæ $v$ valor differentialis
\[=nv.dx\left(-\frac1{dx}d.\frac{WVp}{\sqrt{1+pp}}\right)=-nv.d.\frac{WVp}{\sqrt{1+pp}}.\]
Porro maximi minimive formula $\int\frac{dx\sqrt{1+pp}}{\sqrt v}$ pertinebit ad Casum quintum, eritque $Z=\frac{\sqrt{1+pp}}{\sqrt v}$, \&
\[dZ=-\frac{dv\sqrt{1+pp}}{2v\sqrt v}+\frac{pdp}{\sqrt{v(1+pp)}};\]
ideoque $\pi=v$, \& $L=-\frac{\sqrt{1+pp}}{2v\sqrt v}$, $M=0$, $N=0$, \& $P=\frac p{\sqrt{v(1+pp)}}$. Deinde vero, ob $v=\int dx(g+W\sqrt{1+pp})$, erit $[Z]=g+W\sqrt{1+pp}$, \&
\[d[Z]=Udv\sqrt{1+pp}+\frac{Wpdp}{\sqrt{1+pp}};\]
unde $[L]=U\sqrt{1+pp}$, $[M]=0$, $[N]=0$, \& $[P]=\frac{Wp}{\sqrt{1+pp}}$. Ponatur, si post integrationem fiat $x=a$,
\[-\int e^{\int Udx\sqrt{1+pp}}\frac{dx\sqrt{1+pp}}{2v\sqrt v}=K,\]
\sourcepage{216}sitque
\[e^{-\int Udx\sqrt{1+pp}}\left(K+\int e^{\int Udx\sqrt{1+pp}}\frac{dx\sqrt{1+pp}}{2v\sqrt v}\right)=T:\]
atque erit formulæ $\int\frac{dx\sqrt{1+pp}}{\sqrt v}$ valor differentialis
\[
\begin{aligned}
&=nv.dx\left(\frac1{dx}d\left(\frac p{\sqrt{v(1+pp)}}+\frac{WTp}{\sqrt{1+pp}}\right)\right)\\
&=-nv.d\left(\frac p{\sqrt{v(1+pp)}}+\frac{WTp}{\sqrt{1+pp}}\right).
\end{aligned}
\]
Ex his duobus valoribus differentialibus inventis nascitur pro curva quæsita sequens æquatio,
\[0=\alpha d.\frac{WVp}{\sqrt{1+pp}}+d\left(\frac p{\sqrt{v(1+pp)}}+\frac{WTp}{\sqrt{1+pp}}\right),\]
\& integrando,
\[B=\frac p{\sqrt{v(1+pp)}}+\frac{Wp(\alpha V+T)}{\sqrt{1+pp}}.\]
At est
\[\alpha V+T=e^{-\int Udx\sqrt{1+pp}}\left(\alpha H+K+\int e^{\int Udx\sqrt{1+pp}}\frac{dx\sqrt{1+pp}}{2v\sqrt v}\right).\]
Quod si ergo ponatur $\alpha H+K=C$, erit $C$ constans arbitraria, atque quantitas definita $a$ omnino ex æquatione evanescet; ideoque curva quæsita desideratam proprietatem pro quavis abscissa possidebit. Pro curva quæsita habebitur ergo ista æquatio:
\[
e^{\int Udx\sqrt{1+pp}}\left(\frac{B\sqrt{1+pp}}{Wp}-\frac1{W\sqrt v}\right)=C+\int e^{\int Udx\sqrt{1+pp}}\frac{dx\sqrt{1+pp}}{2v\sqrt v},
\]
\& differentiando
\[
\begin{aligned}
&-\frac{Bdp}{Wp^2\sqrt{1+pp}}-\frac{BUdv\sqrt{1+pp}}{W^2p}+\frac{Udv}{W^2\sqrt v}+\frac{dv}{2Wv\sqrt v}\\
&\quad+\frac{BUdx(1+pp)}{Wp}-\frac{Udx\sqrt{1+pp}}{W\sqrt v}=\frac{dx\sqrt{1+pp}}{2v\sqrt v}.
\end{aligned}
\]
Cum autem sit $dv=gdx+Wdx\sqrt{1+pp}$, habebimus facta substitutione, hanc æquationem
\[\frac{Bdp}{Wp^2\sqrt{1+pp}}=\frac{gdx}{2Wv\sqrt v}+\frac{gUdx}{W^2\sqrt v}-\frac{gBUdx\sqrt{1+pp}}{W^2p},\]
sive istam
\[\frac{2BWdp}{\sqrt{1+pp}}=\frac{gWp^2dx}{v\sqrt v}+\frac{2gUp^2dx}{\sqrt v}-2gBUpdx\sqrt{1+pp}.\]
Multiplicetur hæc æquatio per $dv$, \& in primo termino loco $dv$ scribatur $gdx+\sourcepage{217}Wdx\sqrt{1+pp}$, ac $dW$ loco $Udv$; quo facto, habebitur ista æquatio
\[\frac{2gBdp}{W\sqrt{1+pp}}+2Bdp-\frac{gp^2dv}{Wv\sqrt v}=\frac{2gppdW}{W^2\sqrt v}-\frac{2gBpdW\sqrt{1+pp}}{W^2};\]
quæ divisa per $p^2$ fit integralis; eritque æquatio integrata hæc,
\[2C-\frac{2B}p=\frac{2gB\sqrt{1+pp}}{Wp}=\frac{2g}{W\sqrt v},\]
sive
\[W=\frac{gB\sqrt{v(1+pp)}-gp}{Cp\sqrt v-B\sqrt v}=\frac{dv-gdx}{dx\sqrt{1+pp}}.\]
Unde nascitur æquatio a resistentia $W$ libera hæc,
\[(Cp-B)dv=gCpdx+gBppdx-\frac{gpdx\sqrt{1+pp}}{\sqrt v}.\]
Cum autem $W$ sit functio ipsius $v$ data, ope æquationis
\[W\sqrt v=\frac{gB\sqrt{v(1+pp)}-gp}{Cp-B},\]
dabitur $p$ per $v$; qui valor si in præcedente æquatione substituatur, dabitur $dx$ per $v$ \& $dv$; hincque curva quæsita poterit construi.

\subsection*{Propositio VI. Problema.}
\originalsection{67}Inter omnes curvas proprietate communi $A$ præditas, determinare eam, in qua sit functio quæcunque, cum ipsius illius expressionis $A$, tum alius cujuscunque $B$, maximum vel minimum.

\subsection*{Solutio.}
Sit $dA$ valor differentialis expressionis $A$, atque $dB$ valor differentialis expressionis $B$; habebit functionis illius ipsarum $A$ \& $B$, quam maximum minimumve esse oportet, valor differentialis hujusmodi formam $\alpha dA+\beta dB$; in qua constantes $\alpha$ \& $\beta$ a ratione compositionis qua expressiones $A$ \& $B$ in illa functione inter se permiscentur pendent; ita ut valores obtineant determinatos ab abscissæ quantitate, cui solutionem accommodatam esse oportet pendentes. Quoniam vero expressionis $A$, quæ proprietatem communem complectitur, valor differentialis est
\sourcepage{218}$dA$; hujus multiplum quodcunque $\gamma dA$ addatur ad valorem differentialem $\alpha dA+\beta dB$ expressionis, quæ maximum minimumve esse debet, ac summa $(\alpha+\gamma)dA+\beta dB$ nihilo æqualis posita dabit æquationem pro curva quæsita. Habebitur igitur ista æquatio $(\alpha+\gamma)dA+\beta dB=0$, seu $(\alpha+\gamma)\delta dA+\beta\delta dB=0$; in qua, etiamsi $\alpha$ \& $\beta$ sint quantitates constantes determinatæ, tamen, ob $\gamma$ \& $\delta$ quantitates constantes arbitrarias, coefficientes valorum $dA$ \& $dB$ qui sunt $(\alpha+\gamma)\delta$ \& $\beta\delta$ evadent constantes arbitrariæ magnitudinis. Harum igitur loco si scribantur litteræ $\xi$ \& $\eta$, habebitur pro curva quæsita ista æquatio $\xi dA+\eta dB=0$. Quocirca ad Problema solvendum, expressionum $A$ \& $B$, quarum altera proprietatem communem continet, utriusque autem functio quæcunque maximum minimumve esse debet, singulatim valores differentiales $dA$ \& $dB$ capi oportet, eosque, per quantitates constantes arbitrarias, quasque multiplicatos nihilo æquales poni, quo pacto resultabit ista æquatio $\xi dA+\eta dB=0$, quæ naturam curvæ quæsitæ exprimet. Q. E. I.

\subsection*{Corollarium I.}
\originalsection{68}Natura igitur curvæ satisfacientis tantum ab expressionibus $A$ \& $B$ pendet; neque ratio functionis ipsarum $A$ \& $B$, quæ maximum minimumve esse debet, ullo modo in computo manet; sed quæcunque sit functio, eadem solutio prodibit.

\subsection*{Corollarium II.}
\originalsection{69}Quæcunque itaque ipsarum $A$ \& $B$ functio, inter omnes curvas eadem proprietate $A$ gaudentes, debeat esse maximum vel minimum; solutio perinde se habebit, ac si, inter omnes curvas eadem communi proprietate $A$ gaudentes, ea requiratur, in qua expressio altera $B$ maximum minimumve obtineat valorem.

\subsection*{Corollarium III.}
\originalsection{70}\sourcepage{219}Quod si ergo expressiones $A$ \& $B$ ejusmodi fuerint formulæ, quarum valores differentiales $dA$ \& $dB$ non pendeant a magnitudine abscissæ $x$ cui respondent; quod evenit, si illæ formulæ pertineant ad Casum vel primum vel quartum, secundum nostram enumerationem Capite præcedente §.~7 factam, tum curva inventa pro quacunque abscissa æque satisfaciet.

\subsection*{Corollarium IV.}
\originalsection{71}Eadem Solutio locum habebit si, inter omnes curvas quarum communis sit proprietas functio quæcunque ipsarum $A$ \& $B$, ea requiratur in qua alia quæpiam earundem $A$ \& $B$ functio sit maximum vel minimum. Hoc enim quoque casu pervenitur ad æquationem $\xi dA+\eta dB=0$, in qua $\xi$ \& $\eta$ sint quantitates constantes ad arbitrium accipiendæ.

\subsection*{Exemplum I.}
\originalsection{72}\marginnote{\footnotesize \hyperlink{fig:20}{Fig.~20.}}Inter omnes curvas $aMb$ cum axe $AB$ eandem aream $\int ydx$ continentes, invenire eam in qua sit $\frac{\int yydx}{\int ydx}$ minimum.

Quæstio hæc initur, si inter omnes areas æquales quæ intra ordinatas extremas $Aa$ \& $Bb$ atque basi $AB$ formari possunt, desideretur ea, quæ habeat suum centrum gravitatis in loco infimo positum. Sumpta enim curva quacunque $aMb$, positisque abscissa $AP=x$, applicata $PM=y$, erit portionis $aAPM$ centrum gravitatis a basi $AP$ remotum intervallo $=\frac{\int yydx}{2\int ydx}$; quod adeo fiet minimum, si reddatur hæc expressio $\frac{\int yydx}{\int ydx}$ minima. Habemus ergo binas has formulas $\int ydx$ \& $\int yydx$, quarum valores differentiales sunt $nv.dx.1$ \& $nv.dx.2y$, ex
\sourcepage{220}quibus pro curva quæsita ista colligitur æquatio $\xi+2\eta y=0$; seu $y=c$. Quæstioni igitur satisfacit linea recta $\alpha\beta$ basi $AB$ parallela seu horizontalis, atque parallelogrammum rectangulum $A\alpha\beta B$, præ omnibus aliis figuris ut $AabB$ ejusdem areæ, hac gaudebit prærogativa, ut ejus centrum gravitatis ad basin $AB$ proxime accedat. Quod si ergo $\alpha AB\beta$ concipiatur tanquam vas aqua repletum, si suprema aquæ superficies $\alpha\beta$ sese ad situm horizontalem composuerit, tum aqua habebit suum centrum gravitatis profundius situm, quam si ejus suprema superficies alium quemcunque situm teneret.

\subsection*{Exemplum II.}
\originalsection{73}\marginnote{\footnotesize \hyperlink{fig:14}{Fig.~14.}}Inter omnes curvas ejusdem longitudinis $DAD$, invenire eam quæ habeat suum gravitatis centrum quam profundissime situm, seu in qua sit $\frac{\int xdx\sqrt{1+pp}}{\int dx\sqrt{1+pp}}$ minimum.

Jam intelligitur Solutio hujus Quæstionis datura esse curvam Catenariam; namque secundum leges Staticas catena ex punctis $D$ \& $D$ suspensa ejusmodi induet figuram ut ejus centrum gravitatis maxime descendat. Quamobrem inter omnes figuras, quas catena inducere potest, quæ quidem omnes ejusdem sunt longitudinis, curva Catenaria orietur, si quæratur ea, in qua sit $\frac{\int xdx\sqrt{1+pp}}{\int dx\sqrt{1+pp}}$ minimum; quippe quæ expressio dat distantiam centri gravitatis $G$ ab abscissarum initio $A$. Cum igitur habeantur binæ istæ formulæ $\int dx\sqrt{1+pp}$, \& $\int xdx\sqrt{1+pp}$; quærantur earum valores differentiales, qui erunt, primæ $=-nv.d.\frac p{\sqrt{1+pp}}$, \& alterius $=-nv.d.\frac{xp}{\sqrt{1+pp}}$, ex quibus nascitur pro curva quæsita ista æquatio $cd.\frac p{\sqrt{1+pp}}=d.\frac{xp}{\sqrt{1+pp}}$, \& integrando $\frac{xp}{\sqrt{1+pp}}=\frac{cp}{\sqrt{1+pp}}\sourcepage{221}+b$, seu $x-c=\frac{b\sqrt{1+pp}}p$, \& $dx=\frac{-bdp}{pp\sqrt{1+pp}}$. Hinc ergo fiet $y=\int pdx=-b\int\frac{dp}{p\sqrt{1+pp}}$; ex quibus æquationibus curva construetur, eritque curvæ longitudo $\int dx\sqrt{1+pp}=s=\frac bp+\textit{Const.}=\frac bp+f$. Hinc alia Constructio, definiendis $x$ \& $y$ per $s$ formari poterit: erit nempe $p=\frac b{s-f}$ \& si initium capiatur in $A$, ubi sit $p=\infty$, ponendum est $f=0$, ita ut sit $p=\frac bs$; unde fit $\sqrt{1+pp}=\frac{\sqrt{bb+ss}}s$, \& $dx\sqrt{1+pp}=ds=\frac{dx\sqrt{bb+ss}}s$; hincque $dx=\frac{sds}{\sqrt{bb+ss}}$, \& $x=\sqrt{bb+ss}-b$. Porro erit $dy=pdx=\frac{bds}{\sqrt{bb+ss}}$, atque $y=bl\frac{s+\sqrt{bb+ss}}b$. Æquatio autem inter coordinatas orthogonales $x$ \& $y$ deducetur ex æquatione $x-c=\frac{b\sqrt{1+pp}}p$; quæ si desideretur super axe $AP$, qui est diameter, \& pro initio abscissarum in $A$ sumpto, ubi est $p=\infty$, poni oportet $c=-b$; eritque $(x+b)p=b\sqrt{1+pp}$, hincque $(x+b)^2pp=bb+bbpp$, \& $p=\frac b{\sqrt{xx+2bx}}$; ideoque $dy=\frac{bdx}{\sqrt{xx+2bx}}$, quæ est æquatio pro Catenaria nota.

\subsection*{Exemplum III.}
\originalsection{74}Inter omnes curvas ejusdem longitudinis, determinare eam in qua sit $\frac{\int Sxdx\sqrt{1+pp}}{\int Sdx\sqrt{1+pp}}$ minimum; denotante $S$ functionem quamcunque arcus curvæ $s=\int dx\sqrt{1+pp}$.

In hoc Exemplo continetur inventio curvæ Catenariæ, si catena non fuerit ubique uniformiter crassa, sed cujus crassities arcui\sourcepage{222} $s$ respondens est ut $S$ functio ipsius $s$. Tum enim exprimet $\int Sdx\sqrt{1+pp}$ hujus catenæ pondus, \& $\frac{\int Sxdx\sqrt{1+pp}}{\int Sdx\sqrt{1+pp}}$ altitudinem centri gravitatis supra abscissarum initium; quæ esse debet minima. Principio quidem hic casus in Problemate præcedente non contineri videtur, quia formula arcum exprimens ipsa $\int dx\sqrt{1+pp}$ non inest in maximi minimive expressione $\frac{\int Sxdx\sqrt{1+pp}}{\int Sdx\sqrt{1+pp}}$, quippe quæ est functio duarum aliarum formularum integralium. At cum sit $S$ functio arcus curvæ $s$, atque $ds=dx\sqrt{1+pp}$, erit $\int Sdx\sqrt{1+pp}=\int Sds$, ideoque functio ipsius $s$: ex quo expressio $\frac{\int Sxdx\sqrt{1+pp}}{\int Sdx\sqrt{1+pp}}$, erit functio formularum $\int dx\sqrt{1+pp}$ \& $\int Sxdx\sqrt{1+pp}$, quarum illa proprietatem communem continet. Idem igitur est ac si quærere deberemus inter omnes curvas æque longas eam in qua sit $\int Sxdx\sqrt{1+pp}$ minimum. Cum jam $S$ sit functio ipsius $s=\int dx\sqrt{1+pp}$, pertinebit hæc Quæstio ad Propositionem præcedentem, eumque casum qui §.~60 est pertractatus. Scilicet erit $Z=Sx\sqrt{1+pp}$; unde, si ponamus $dS=Tds$, fiet
\[dZ=xTds\sqrt{1+pp}+Sdx\sqrt{1+pp}+\frac{Sxpdx}{\sqrt{1+pp}},\]
ita ut sit $L=xT\sqrt{1+pp}$; $M=S\sqrt{1+pp}$, $N=0$ \& $P=\frac{Sxp}{\sqrt{1+pp}}$. Jam ob $N=0$, obtinemus ex eodem loco citato statim hanc æquationem
\[A+\frac{(C-\int xTdx\sqrt{1+pp})p}{\sqrt{1+pp}}=\frac{-Sxp}{\sqrt{1+pp}},\]
seu
\[\frac{A\sqrt{1+pp}}p+C-\int xTdx\sqrt{1+pp}+Sx=0.\]
At est $Tdx\sqrt{1+pp}=Tds=dS$; unde habetur $\frac{A\sqrt{1+pp}}p+C+Sx-\int xdS=0$, ubi $A$ \& $C$ sunt quantitates arbitrariæ. Differentietur hæc æquatio, fietque $\frac{-Adp}{pp\sqrt{1+pp}}+Sdx=0$, seu $Sdx\sqrt{1+pp}=-\frac{cdp}{pp}=Sds$. Quare cum sit $S$ functio
\sourcepage{223}ipsius $s$, integretur $Sds$, eritque integrale, quod sit $=R$, pondus catenæ longitudini $s$ respondens. Fiet ergo integrando $\frac cp=R+C$; \&, si initium curvæ capere placeat in loco $A$, ubi curvæ tangens est horizontalis, erit $C=0$, atque $p=\frac cR$. Hinc ergo porro erit $\sqrt{1+pp}=\frac{\sqrt{cc+RR}}R=\frac{ds}{dx}$; ideoque $dx=\frac{Rds}{\sqrt{cc+RR}}$, atque $dy=\frac{cds}{\sqrt{cc+RR}}$, ex quibus æquationibus curva ita poterit construi, ut statim ad quamvis catenæ longitudinem tam abscissa quam applicata respondens definiatur. Manifestum autem est casu quo $R=s$, hoc est quo catena ponitur uniformis crassitiei, tum prodire Catenariam curvam ordinariam.

\subsection*{Scholion.}
\originalsection{75}Nisi hujus Exempli convenientia, tam cum ista Propositione quam cum præcedente, esset observata, tum Solutio quidem per regulam generalem absolvi potuisset: verum tamen multo prolixior evasisset. Quo autem nihilominus Methodi generalis usus clarius ob oculos ponatur, idem hoc Exemplum secundum generalia præcepta resolvere visum est. Quæratur igitur inter omnes curvas ejusdem longitudinis $s=\int dx\sqrt{1+pp}$, ea quæ habeat valorem expressionis hujus $\frac{\int Sxdx\sqrt{1+pp}}{\int Sdx\sqrt{1+pp}}$ maximum vel minimum; existente $S$ functione quacunque arcus curvæ $s$. Et quoniam nondum suspicari licet considerationem datæ abscissæ, a qua valor differentialis expressionis $\frac{\int Sxdx\sqrt{1+pp}}{\int Sdx\sqrt{1+pp}}$ pendet, ex calculo esse egressuram; ponamus huic Quæstioni tantum pro data abscissæ longitudine $x=a$ satisfieri oportere. Ab hac longitudine quidem formulæ communem proprietatem continentis $\int dx\sqrt{1+pp}$ valor differentialis non pendet,
\sourcepage{224}quippe qui constanter est $=-nv.d.\frac p{\sqrt{1+pp}}$; at in maximi minimive expressione $\frac{\int Sxdx\sqrt{1+pp}}{\int Sdx\sqrt{1+pp}}$ ponamus, casu quo $x=a$, fore $\int Sxdx\sqrt{1+pp}=A$ \& $\int Sdx\sqrt{1+pp}=B$: illius vero numeratoris $\int Sxdx\sqrt{1+pp}$ valorem differentialem esse $=dA$, denominatoris vero $\int Sdx\sqrt{1+pp}$ valorem differentialem esse $=dB$. Hinc igitur maximi minimive expressionis, quæ, casu $x=a$, fit $=\frac AB$, valor differentialis erit $\frac{BdA-AdB}{BB}$, qui multiplo cuicunque formulæ communis valoris differentialis $-nv.d.\frac p{\sqrt{1+pp}}$ æqualis positus dabit æquationem pro curva quæsita. Jam ad valores differentiales $dA$ \& $dB$ inveniendos, consideremus primum formulam $\int Sdx\sqrt{1+pp}$, quæ secundum enumerationem §.~7 Cap. præced. factam, pertinet ad Casum secundum: quo erit $Z=S\sqrt{1+pp}$, \& posito $dS=Tds$, erit
\[dZ=Tds\sqrt{1+pp}+\frac{Spdp}{\sqrt{1+pp}}.\]
Comparatione ergo facta, erit $\pi=s$; $L=T\sqrt{1+pp}$, $M=0$, $N=0$, \& $P=\frac{Sp}{\sqrt{1+pp}}$: tum vero ob $\pi=s=\int dx\sqrt{1+pp}$, erit $[Z]=\sqrt{1+pp}$ \& $[M]=0$, $[N]=0$, \& $[P]=\frac p{\sqrt{1+pp}}$. Jam sumatur integrale $\int Ldx=\int Tdx\sqrt{1+pp}=\int Tds=S$, cujus valor, casu $x=a$, fiat $=G$, eritque $V=G-S$. Quamobrem habebitur formulæ $\int Sdx\sqrt{1+pp}$ valor differentialis
\[dB=-nv.d\left(\frac{Sp}{\sqrt{1+pp}}+\frac{p(G-S)}{\sqrt{1+pp}}\right)=-nv.d.\frac{Gp}{\sqrt{1+pp}}.\]
Altera porro formula $\int Sxdx\sqrt{1+pp}$ pariter in eodem Casu secundo comprehenditur, eritque $Z=Sx\sqrt{1+pp}$ \&
\[dZ=Txds\sqrt{1+pp}+Sdx\sqrt{1+pp}+\frac{Sxpdp}{\sqrt{1+pp}},\]
unde fit $\pi=s$; $L=Tx\sqrt{1+pp}$; $M=S\sqrt{1+pp}$;
\sourcepage{225}$N=0$ \& $P=\frac{Sxp}{\sqrt{1+pp}}$. Deinde, ob $\pi=s=\int dx\sqrt{1+pp}$, erit, ut ante, $[Z]=\sqrt{1+pp}$, $[M]=0$, $[N]=0$, \& $[P]=\frac p{\sqrt{1+pp}}$. Nunc sumatur integrale $\int Ldx=\int Txdx\sqrt{1+pp}=\int Txds=\int xdS$, cujus valor, posito $x=a$, sit $=H$; erit $V=H-\int xdS$, hincque prodibit istius formulæ valor differentialis
\[\begin{aligned}dA&=-nv.d\left(\frac{Sxp+p(H-\int xdS)}{\sqrt{1+pp}}\right)\\&=-nv.d\left(\frac{p(H+\int Sdx)}{\sqrt{1+pp}}\right).\end{aligned}\]
Inventis ergo valoribus $dA$ \& $dB$, æquatio pro curva quæsita erit
\[\alpha B^2d.\frac p{\sqrt{1+pp}}-\beta Ad.\frac{Gp}{\sqrt{1+pp}}+\beta Bd.\frac{p(H+\int Sdx)}{\sqrt{1+pp}}=0,\]
\& integrando
\[\frac{\alpha B^2p-\beta AGp+\beta BHp+\beta Bp\int Sdx}{\sqrt{1+pp}}=C;\]
in qua $\alpha$, $\beta$, \& $C$ sunt constantes arbitrariæ, \& $G$ \& $H$ constantes determinatæ. Quod si ergo ponatur $\frac{\alpha B}{\beta}+\frac{AG}{B}+H=b$ \& $\frac C{\beta.B}=c$, erunt $b$ \& $c$ constantes arbitrariæ, atque constantes determinatæ $G$ \& $H$ a definito abscissæ valore $x=a$ pendentes omnino ex æquatione evanescent; ita ut Curva inventa pro quavis abscissa gavisura sit desiderata proprietate: ejusque æquatio erit hæc
\[c=\frac{bp+p\int Sdx}{\sqrt{1+pp}},\qquad\text{seu}\qquad\frac{c\sqrt{1+pp}}p=b+\int Sdx;\]
quæ differentiata dabit $Sdx=\frac{-cdp}{pp\sqrt{1+pp}}$, seu $Sdx\sqrt{1+pp}=Sds=\frac{-cdp}{pp}$. Ponatur, ut supra, $\int Sds=R$, ita ut $R$ pondus longitudinis catenæ $s$ repræsentet, erit $R=\frac cp+\text{Const.}$ quæ est ipsa æquatio, quam præcedenti Methodo elicuimus. Ex hac itaque solutione intelligitur, quemadmodum per Methodum generalem hujusmodi Quæstiones resolvi possint, si proprietas communis non ingrediatur in maximi minimive expressionem; quod ut clarius intelligatur unum adhuc hujusmodi Exemplum apposuisse sufficiet.

\sourcepage{226}\subsection*{Exemplum IV.}
\originalsection{76}\marginnote{\footnotesize \hyperlink{fig:14}{Fig.~14.}}Inter omnes curvas ejusdem longitudinis $DAD$ datæ abscissæ $AC=a$ respondentes, eam definire quæ comprehendat aream $DAD$, cujus centrum gravitatis $G$ sit vel altissime vel profundissime positum, seu in qua sit $\frac{\int yxdx}{\int ydx}$ maximum vel minimum.

Proprietas igitur communis est $\int dx\sqrt{1+pp}$, cujus valor differentialis cuicunque abscissæ $x$ respondens est $=-nv.d.\frac p{\sqrt{1+pp}}$. Maximi autem minimive expressionis $\frac{\int yxdx}{\int ydx}$ valor differentialis pendebit a præscripta abscissæ longitudine $x=a$; qui ut inveniatur; casu quo $x=a$, fiat $\int yxdx=A$, hujusque formulæ valor differentialis sit $=dA$, qui per Regulas supra datas invenitur $=nv.dx.x=nv.xdx$. Porro, eodem casu $x=a$, abeat altera formula $\int ydx$ in $B$, sitque ejus valor differentialis $=dB$, qui per Regulas datas reperitur $=nv.dx$; ita ut sit $dA=nv.xdx$ \& $dB=nv.dx$. Ex his, maximi minimive expressionis $\frac{\int yxdx}{\int ydx}$, quæ, casu $x=a$, abit in $\frac AB$, valor differentialis erit $\frac{BdA-AdB}{BB}=\frac{nv(Bxdx-Adx)}{BB}$, qui multiplo valoris differentialis $-nv.d.\frac p{\sqrt{1+pp}}$, qui ex proprietate communi prodiit, æqualis positus dabit pro curva quæsita istam æquationem
\[\alpha d.\frac p{\sqrt{1+pp}}=\frac{Bxdx-Adx}{BB}.\]
Sit $\frac AB=h$, erit $h$ quantitas constans determinata, quam præbet formula $\frac{\int yxdx}{\int ydx}$, si ponatur $x=a$, \& $\alpha B$ ponatur $=cc$, erit $cc$ quantitas arbitraria. Hinc habebitur ista pro curva æquatio
\[ccd.\frac p{\sqrt{1+pp}}=xdx-hdx,\]
\sourcepage{227}quæ integrata dat $\frac{2ccp}{\sqrt{1+pp}}=xx-2hx+bb$; ergo $4c^4pp=(xx-2hx+bb)^2(1+pp)$, atque
\[p=\frac{xx-2hx+bb}{\sqrt{4c^4-(xx-2hx+bb)^2}}=\frac{dy}{dx}.\]
Quocirca erit
\[y=\int\frac{(xx-2hx+bb)dx}{\sqrt{4c^4-(xx-2hx+bb)^2}},\]
ubi constantem $bb$, pro arbitrio, sive affirmativam, sive negativam accipere licet. Hæc autem curva Quæstioni satisfacit tantum casu, quo $x=a$; atque ut satisfaciat litteræ $h$ is tribui debet valor quem, casu $x=a$, recipiet expressio $\frac{\int yxdx}{\int ydx}$, ex quo valor $h$ determinabitur. Cæterum notari convenit hanc curvam esse eam quæ vulgo sub nomine Elasticæ est cognita.
